\documentclass[11pt,a4paper]{article}
\usepackage{fontspec}
\usepackage{amsmath,amssymb}
\usepackage{booktabs}
\usepackage{longtable}
\usepackage{array}
\usepackage{enumitem}
\usepackage{xcolor}
\usepackage[margin=2.2cm]{geometry}
\usepackage{fancyhdr}
\usepackage[colorlinks=true,linkcolor=black,citecolor=black,urlcolor=blue]{hyperref}

\setmainfont{DejaVu Sans}
\setmonofont{NotoSansMono-Regular.ttf}[Scale=0.94]
\definecolor{okgreen}{RGB}{0,110,60}

\pagestyle{fancy}
\fancyhf{}
\lhead{\small Leptonic compactness and the generalized $n$-dimensional law}
\rhead{\small\thepage}
\renewcommand{\headrulewidth}{0.3pt}

\setlength{\parskip}{3pt}
\setlength{\parindent}{0pt}

\title{}
\date{}
\begin{document}
\sloppy

\begin{center}
{\LARGE\bfseries Leptonic Compactness and the Generalized\\[2pt]
$n$-Dimensional Velocity Law}\\[6pt]
{\large A formal report on the \texttt{TheoreticalFramework} development}\\[10pt]
{\small Generated 2026-11-05 13:39 UTC from the Lean sources by \texttt{paper/tools/build\_paper.py}}
\end{center}

\vspace{4pt}
\begin{center}
\begin{minipage}{0.92\textwidth}
\small
\textbf{Status:} \textbf{\textcolor{okgreen}{all checks passed}}.\\
\textbf{Lean:} \texttt{4.33.1}; \textbf{Mathlib:} prebuilt at \texttt{/opt/mathlib}.\\
\textbf{Sources:} 5 modules, 97 theorems, 42 definitions,
3 structures.\\
\textbf{Axioms used across the development:} \texttt{Classical.choice}, \texttt{Quot.sound}, \texttt{propext}.
\end{minipage}
\end{center}

\vspace{6pt}
\noindent\rule{\textwidth}{0.4pt}
\section*{Summary}
\addcontentsline{toc}{section}{Summary}

Part of the muon's anomalous magnetic moment is attributed, in the framework
under study, to an internal geometric compactness of massive leptons with
effective radius $r_\ell = \kappa/m_\ell$.  Under the framework's plate ansatz
($S$ proportional to $m^2$, separation $d = r_\ell$) the Casimir energy of the
internal geometry is quintic in the mass, so the amplification between two
leptons is the fifth power of their mass ratio:
$E_{\mathrm{C},\mu}/E_{\mathrm{C},e} \approx (207)^5 \approx 4\times10^{11}$
and $E_{\mathrm{C},\tau}/E_{\mathrm{C},e} \approx (3477)^5 \approx 5.1\times10^{17}$.
Within the generalized $n$-dimensional kinematics with $n = 5$, the same quintic
dependence fixes the ejection-velocity law
$v_f^2 = v_2^2\,(m_\ell/m_0)^{n-2}$, hence amplification factors
$(207)^3 \approx 8.9\times10^{6}$ and $(3477)^3 \approx 4.2\times10^{10}$.

This report does three things.  First, it records the derivation of every
scaling law as a Lean theorem, so that the fifth-power and cube appear as
consequences rather than assertions, and so that the parameter-freedom of the
ratios is visible.  Second, it recomputes every physical constant independently
in Python and compares it with the interval Lean states --- including the three
places where the source text's arithmetic does not survive that comparison.
Third, it separates what the development proves from what the framework assumes
and from what would need experiment.


\section{Scope, and what is and is not proved}
\label{sec:scope}

The development is deliberately layered, and the layers should not be confused.

\paragraph{(A) Proved mathematics.}
The scaling laws, the amplification ratios, the exact arithmetic of the ejecta,
the algebra of the regime structure, the continuity of the velocity map, and the
thrust--power identity.  These are Lean theorems with no hypotheses beyond
positivity of the masses and ratios, and they hold whatever one thinks of the
physics.

\paragraph{(B) Framework postulates.}
Two assumptions are \emph{recorded} rather than derived, and the formalization
makes that visible instead of hiding it.  The geometric ansatz
($S$ proportional to $m^2$ with $d = r_\ell$, and $r_\ell = \kappa/m_\ell$) is a
structure (\texttt{CasimirAnsatz}) carrying a free constant $\kappa$ and a free
area coefficient; the theorem that matters is that the \emph{ratio} does not
depend on either of them
(\texttt{CasimirAnsatz.energy\_ratio\_parameter\_free}).  The velocity law
$v_f^2 = v_2^2\,(m_\ell/m_0)^{n-2}$ is likewise a postulate of the framework's
generalized kinematics, and the development proves its consequences
(\texttt{velocity\_law}), not its truth.

\paragraph{(C) Empirical claims.}
That any of this describes the muon's actual internal structure, that muon--muon
collisions eject particles at these speeds, and that a thruster can be built, are
questions for experiment.  The development is explicit about the most important
of these: \S\ref{sec:thrust} proves that the geometry fixes the exhaust velocity
and the thrust-to-power ratio but says nothing about absolute thrust, and the
claim that is left open is recorded as a proposition
(\texttt{ThrustOptimizationConjecture}) that is never asserted.

\paragraph{A tension internal to the source.}
The framework attributes part of $a_\mu$ to a term scaling as $m^5$, whereas the
measured generation-dependence of the anomalous moment goes as roughly
$(m_\mu/m_e)^2 \approx 4.3\times10^{4}$, a factor $\approx 10^{7}$ smaller than
$(207)^5$.  The attached notes say so themselves.  Nothing in this report
resolves that: the formalization establishes that the framework's own arithmetic
is correct, which is a different statement from the framework being right.  A
Lean proof of a scaling law is not evidence for the postulate that produces it.


\section{Core derivation: compactness and the fifth power}
\label{sec:core}

Write the effective radius as $r_\ell = \kappa/m_\ell$ and the Casimir energy of
the internal cavity as $E_\mathrm{C} \propto -S/d^3$ with $d = r_\ell$.  Three
theorems carry the derivation.

\begin{enumerate}[leftmargin=*]
\item \texttt{casimirEnergy\_eq\_quintic}: for $\kappa \neq 0$ and $m \neq 0$,
  \[
    E_\mathrm{C}(\kappa,m) \;=\; \frac{m^2}{(\kappa/m)^3} \;=\; \frac{m^5}{\kappa^3}.
  \]
\item \texttt{casimirEnergy\_ratio}: for non-zero masses,
  \[
    \frac{E_\mathrm{C}(\kappa,m_1)}{E_\mathrm{C}(\kappa,m_2)} \;=\; \left(\frac{m_1}{m_2}\right)^{5}.
  \]
\item \texttt{casimirEnergy\_ratio\_radii}: the same number read geometrically,
  \[
    \frac{E_\mathrm{C}(\kappa,m_1)}{E_\mathrm{C}(\kappa,m_2)} \;=\; \left(\frac{r_2}{r_1}\right)^{5},
  \]
  so ``compactness'' and ``vacuum coupling'' are one quantity, not two.
\end{enumerate}

The constant $\kappa$ cancels in the ratio.  That is the whole reason the
framework can issue a number at all, and the formalization isolates it: in
\texttt{CasimirAnsatz} the ratio is proved independent of $\kappa$, of the
area coefficient, and of the overall Casimir prefactor, and the sign convention
(the physical energy is negative, i.e.\ attractive) is shown not to affect it
(\texttt{signedEnergy\_ratio}).

The effective radius is strictly decreasing in the mass
(\texttt{effRadius\_strictAnti}), which is the compactness statement in the form
the rest of the development uses.


\section{Casimir amplification: the numbers}
\label{sec:casimir}

With $m_\mu/m_e = 207$ and $m_\tau/m_e = 3477$ --- both checked against the
tabulated masses to better than a quarter and a half unit respectively
(\texttt{massRatioMuon\_consistent}, \texttt{massRatioTauon\_consistent}) ---
the fifth-power law gives the following.  ``Lean'' is the interval the theorem
\texttt{\dots\_order} states; ``recomputed'' is this script's independent
evaluation; ``published'' is the source text.

\begin{center}\small
\begin{tabular}{@{}lccc@{}}
\toprule
Quantity & Lean interval & Recomputed & Published \\
\midrule
$E_{\mathrm{C},\mu}/E_{\mathrm{C},e}$ & $[3.8 \times 10^{11},\ 4 \times 10^{11}]$ & $3.800596 \times 10^{11}$ & $4\times10^{11}$ \\
$E_{\mathrm{C},\tau}/E_{\mathrm{C},e}$ & $[5.08 \times 10^{17},\ 5.1 \times 10^{17}]$ & $5.081869 \times 10^{17}$ & $5.1\times10^{17}$ \\
$E_{\mathrm{C},\tau}/E_{\mathrm{C},\mu}$ & $[1.3 \times 10^{6},\ 1.34 \times 10^{6}]$ & $1.337124 \times 10^{6}$ & $1.3\times10^{6}$ \\

\bottomrule
\end{tabular}
\end{center}

Two further structural facts are proved rather than computed.  The amplification
is monotone in the mass ratio
(\texttt{casimir\_amplification\_hierarchy}: $1 < (207)^5 < (3477)^5$), so the
lepton generation hierarchy is a strict hierarchy of vacuum couplings.  And the
equivalent magnetic field $B$, defined by $B^2 \propto E_\mathrm{C}/V$, carries
exponent $3$ under the source's volume ansatz $V$ proportional to $S d^3$, and
exponent $2$ under the literal plate volume $V = S d$.  Both are proved
(\texttt{field\_amplification\_source}, \texttt{field\_amplification\_plate}).
The exponent $3$ is the one the framework uses, and it is the same $3 = n-2$
that appears in the velocity law --- which is worth stating precisely, because
the coincidence is the framework's only argument for $n = 5$ beyond the quintic
mass dependence itself.


\section{The generalized $n$-dimensional framework}
\label{sec:dimensions}

The generalized dimension is not tuned: it is read off the quintic mass
dependence of the anomalous-moment term.  In the development,

\begin{verbatim}
casimirMassExponent   = 5
generalizedDimension  = casimirMassExponent
velocityExponent      = generalizedDimension - 2  = 3
\end{verbatim}

\noindent so that \texttt{generalizedDimension\_eq\_five} and
\texttt{velocityExponent\_eq\_three} are definitional, and the exponent of the
velocity law is a consequence of the Casimir exponent rather than a choice made
independently of it.

The link to $g-2$ is made explicit rather than left in prose.  The framework's
decomposition $a_\ell = a_{\mathrm{QED}} + a_{\mathrm{EW+had}} + a_{\mathrm{str}}$
is the definition \texttt{anomalousMoment}, its structural part is
\texttt{structuralTerm} $= C\,m^{\,n}$, and \texttt{structuralTerm\_quintic}
records that this is $C\,m^5$ while
\texttt{structuralTerm\_ratio} records that the structural term amplifies by
$(m_1/m_2)^5$ --- the same number as the Casimir amplification.  That identity of
exponents is the whole basis for reading $n = 5$ off the anomalous moment, so it
is worth having as a theorem rather than a remark.

The law is stated for arbitrary $n$ and then specialized.  With $r = m_\ell/m_0$,
\[
  v_f \;=\; \sqrt{v_2^2\, r^{\,n-2}},
  \qquad\text{so}\qquad
  v_f^2 \;=\; v_2^2\, r^{\,n-2}
\]
(\texttt{velocity\_law}), and for $v_2 \ge 0$ this is the paper's multiplicative
form $v_f = v_2\sqrt{r^{\,n-2}}$ (\texttt{generalisedVelocity\_eq\_mul}).  At
$n = 5$ both become the cube (\texttt{velocity\_law\_five}).  The map is
monotone in $r$ and in $v_2$, vanishes at $v_2 = 0$, and is continuous in each
variable.

The framework's generalized units enter only as exponent bookkeeping, and that
bookkeeping is proved exactly: with $[v_n] = \mathrm{m}^{2/n}/\mathrm{s}$ the mass
exponent balances as $n \cdot (2/n) = 2$
(\texttt{generalised\_unit\_exponents}) and the time exponent as
$(n-2)+2 = n$ (\texttt{time\_exponent\_balance}).  The substitution rule
$v_n (q\,dt_0)^{n-2} = v_2^2 q^{\,n-2}$ that links the two regimes is
\texttt{scale\_substitution}.


\section{Velocity scaling and amplification}
\label{sec:velocity}

The amplification factors are exact integers:
$207^3 = 8\,869\,743$ and $3477^3 = 42\,035\,292\,333$
(\texttt{muonAmplification\_eq}, \texttt{tauonAmplification\_eq}), i.e.\
$\approx 8.9\times10^{6}$ and $\approx 4.2\times10^{10}$.  Their ratio is
$(3477/207)^3 \approx 4.74\times10^{3}$.

Because $v_f^2$ is an exact rational, every velocity is bounded without any
rounding step: the development computes $v_f^2$ exactly and closes the interval
with \texttt{generalisedVelocity\_bounds}.

\begin{center}\small
\begin{tabular}{@{}lccc@{}}
\toprule
Quantity & Lean interval (m/s unless stated) & Recomputed & Published \\
\midrule
muonic ceiling $v_f$ & $[2.0682 \times 10^{7},\ 2.06821 \times 10^{7}]$ & $2.068202 \times 10^{7}$ & $\approx 20\,700$ km/s \\
muonic ceiling (km/s) & $[20{,}680,\ 20{,}700]$ & $20{,}682$ & --- \\
muonic $I_{\mathrm{sp}}$ (s) & $[2.108 \times 10^{6},\ 2.109 \times 10^{6}]$ & $2.108979 \times 10^{6}$ & $\approx 2\times10^{6}$ \\
$I_{\mathrm{sp}}/I_{\mathrm{sp,Raptor}}$ & $[5{,}800,\ 5{,}820]$ & $5{,}810$ & $\approx 6000$ \\
tauonic ceiling $v_f$ & $[4.2403 \times 10^{12},\ 4.2404 \times 10^{12}]$ & $4.240334 \times 10^{12}$ & $1.34\times10^{12}$ \\
tauonic at entry & $[1.4237 \times 10^{9},\ 1.4239 \times 10^{9}]$ & $1.423785 \times 10^{9}$ & $1.42\times10^{9}$ \\
tauonic at median & $[2.8475 \times 10^{10},\ 2.8476 \times 10^{10}]$ & $2.847571 \times 10^{10}$ & $9.0\times10^{10}$ \\
\bottomrule
\end{tabular}
\end{center}

\noindent The ceiling is at the muonic collision maximum $v_2 = 25\,000$ km/h
$= 62500/9 \approx 6\,944.4$ m/s (\texttt{collisionVelocityMuonMax\_eq}).

\subsection*{The law over the whole input range}

Values computed from the verified law $v_f^2 = v_2^2\, r^{3}$; only the endpoints
are bounded inside Lean, the rest are consequences.

\begin{center}\small
\begin{tabular}{@{}rccc@{}}
\toprule
$v_2$ (km/h) & $v_2$ (m/s) & $v_f$ muonic (m/s) & $v_f$ tauonic (m/s) \\
\midrule
0 & $0$ & $0$ & $0$ \\
2,500 & $694$ & $2.068202 \times 10^{6}$ & $1.423785 \times 10^{8}$ \\
5,000 & $1{,}389$ & $4.136405 \times 10^{6}$ & $2.847571 \times 10^{8}$ \\
10,000 & $2{,}778$ & $8.272809 \times 10^{6}$ & $5.695142 \times 10^{8}$ \\
25,000 & $6{,}944$ & $2.068202 \times 10^{7}$ & $1.423785 \times 10^{9}$ \\
50,000 & $13{,}889$ & $4.136405 \times 10^{7}$ & $2.847571 \times 10^{9}$ \\
250,000 & $69{,}444$ & $2.068202 \times 10^{8}$ & $1.423785 \times 10^{10}$ \\
1,000,000 & $2.777778 \times 10^{5}$ & $8.272809 \times 10^{8}$ & $5.695142 \times 10^{10}$ \\
\bottomrule
\end{tabular}
\end{center}

\noindent Rows above $25\,000$ km/h lie outside the muonic input window and are
tabulated only to show the law's behaviour over three decades; they are not
outputs of a regime.  The muonic column above $\approx 10^{5}$ km/h is already
superluminal, which is one reason the muonic window stops where it does.

\subsection*{Where the source text diverges}

Three published figures were checked against the exact law.  Two do not survive,
and the pattern is the same in both: the source computes $v_f^2$ correctly and
then takes its square root with the decimal exponent misplaced, which produces a
factor $\sqrt{10} \approx 3.16$ on $v_f$ whenever the exponent of $v_f^2$ is odd.

\begin{center}\small
\begin{tabular}{@{}lcccl@{}}
\toprule
Quantity & Published & Verified & Ratio & Diagnosis \\
\midrule
tauonic ceiling & $1.34\times10^{12}$ & $4.2403\times10^{12}$ & $3.1644$
  & $\sqrt{10}$ slip \\
tauonic at median & $9.0\times10^{10}$ & $2.8476\times10^{10}$ & $0.3164$
  & same slip, other direction \\
tauonic at entry & $1.42\times10^{9}$ & $1.4238\times10^{9}$ & $1.003$
  & correct (even exponent) \\
\bottomrule
\end{tabular}
\end{center}

The slip is not asserted but localised: \texttt{source\_tauon\_max\_velocity\_underestimates}
proves that the published figure squared is a factor in $[10.01, 10.02]$ below
the verified one, and \texttt{source\_tauon\_max\_velocity\_factor} that the
figure itself is low by a factor in $[3.164, 3.165]$.  The source's own
intermediate, $v_f^2 \approx 1.80\times10^{25}$ m²/s², is right; only its root is
wrong.  Both values are kept in the development
(\texttt{sourceTauonMaxVelocity}) so that nothing is silently corrected.

\subsection*{The speed limit}

The muonic ceiling is subluminal (\texttt{muonMaxEjection\_subluminal}:
$2.07\times10^{7} < c$), and so is the muonic collision window.  The tauonic
ceiling is not (\texttt{tauonMaxEjection\_superluminal}:
$c < 4.24\times10^{12}$ m/s).  The tauonic \emph{input} window is still
subluminal (\texttt{tauonicInputMax\_subluminal}), so the framework's rejection
of the Lorentz factor is required only to interpret the output.  This is the
sharpest formal statement of the framework's distance from standard kinematics,
and it is worth being explicit: a superluminal exhaust is not a technical
difficulty to be engineered around, it is a consequence of the law as written.


\section{The two regimes, and what continuity actually says}
\label{sec:regimes}

Both regimes are instances of one structure, \texttt{PropulsiveRegime}, carrying
a mass ratio, an input window and the law-induced output.  Everything proved of
the structure therefore applies to both.

\begin{center}\small
\begin{tabular}{@{}lcc@{}}
\toprule
 & muonic & tauonic \\
\midrule
Colliding leptons & $\mu^-/\mu^+$ & $\tau^-/\tau^+$ \\
Mass ratio $m/m_e$ & 207 & 3477 \\
Amplification $(m/m_e)^3$ & $8.87\times10^{6}$ & $4.20\times10^{10}$ \\
Casimir amplification & $3.80\times10^{11}$ & $5.08\times10^{17}$ \\
Input window $v_2$ & $]0,\ 25\,000\ \mathrm{km/h}]$ & $]25\,000\ \mathrm{km/h},\ v_f^{\max,\mu}]$ \\
Output ceiling & $20\,682\ \mathrm{km/s}$ & $4.2403\times10^{12}\ \mathrm{m/s}$ \\
\bottomrule
\end{tabular}
\end{center}

\paragraph{Continuity, stated correctly.}
The source text says that ``the entry velocity of the tauonic domain coincides
exactly with the maximum output velocity of the muonic domain''.  Those are two
different coincidences, and the formalization separates them, because conflating
them is how a boundary argument goes wrong:

\begin{enumerate}[leftmargin=*]
\item \textbf{Adjacency} (\texttt{windows\_adjacent}): the tauonic window begins
  at $25\,000$ km/h, which is the muonic \emph{input} ceiling.
\item \textbf{The chain closes} (\texttt{chain\_closes}): the tauonic window's
  \emph{upper} end is the muonic regime's maximum \emph{output}, $20\,682$ km/s.
  So the second regime's input ceiling is the first regime's output ceiling ---
  true of the upper end, not the entry.
\end{enumerate}

Together with \texttt{regimes\_tile} (the two windows cover one interval, with no
gap and no overlap) and \texttt{tauonic\_window\_nested\_in\_output} (the tauonic
window sits inside the muonic output range), this is the exact content of the
nested-regime claim.  \texttt{regime\_continuity} collects the four facts and the
continuity of the velocity map in $v_2$ in one statement.

The boundary itself is a named object, \texttt{domainBoundary}, with
\texttt{domainBoundary\_eq\_muonic\_output} (it *is* the muonic output ceiling)
and \texttt{domainBoundary\_km\_per\_s} (it lies in $[20\,680, 20\,700]$ km/s,
i.e.\ the source's $20\,700$ to the precision of the print).  It is subluminal
(\texttt{domainBoundary\_subluminal}).

Taking the unrounded value $v_f^{\max,\mu}$ as the tauonic input ceiling rather
than the printed $20\,700$ km/s is what makes the match exact rather than
approximate: the theorem \texttt{chain\_closes} is then true by construction,
whereas with the rounded figure it would hold only to three significant digits.
That choice is recorded in the definition of \texttt{tauonicInputMax} and
justified in its docstring.

\paragraph{Structural algebra.}
The amplification $r \mapsto r^{\,n-2}$ is a multiplicative character
(\texttt{amplification\_mul}), so the tauonic factor factors through the muonic
one (\texttt{tauon\_amplification\_factors}):
$3477^3 = 207^3 \cdot (3477/207)^3$.  This is why the regimes compose: a
two-stage cascade multiplies, and the ordering of the regimes is forced by the
lepton mass hierarchy rather than chosen (\texttt{regimes\_strictly\_ordered}).


\section{Thrust: what the framework fixes, and the question it leaves open}
\label{sec:thrust}

The kinematics fix the exhaust velocity.  Momentum and energy balance then fix
the thrust-to-power ratio exactly, and nothing else:

\[
  T = \dot m\, v_e, \qquad P = \tfrac12 \dot m\, v_e^2
  \quad\Longrightarrow\quad
  \frac{T}{P} = \frac{2}{v_e}
\]

\noindent proved as \texttt{thrust\_to\_power} for an arbitrary
\texttt{ThrusterDesign}, and as \texttt{thrust\_eq\_power}: $T = 2P/v_e$.

This has a consequence the source text does not mention and which the
development makes unavoidable.  At the muonic exhaust velocity,
$T/P \approx 9.67\times10^{-8}$ N/W
(\texttt{muonic\_thrust\_to\_power\_bounds}), which is a factor
$\approx 5.8\times10^{3}$ \emph{worse} than a Raptor engine
(\texttt{muonic\_thrust\_to\_power\_shortfall}).  The specific-impulse gain and
the thrust-per-watt loss are the same number, because both are $v_e$: raising
$I_\mathrm{sp}$ from $363$ s to $2.1\times10^{6}$ s raises it by $5.8\times10^{3}$
and lowers thrust per watt by $5.8\times10^{3}$.  Any muonic thruster trades one
for the other; there is no regime in which the law gives both.

What the framework does \emph{not} fix is absolute thrust.  This is proved, not
asserted: \texttt{thrust\_not\_determined} shows that for any positive thrust
$T$ there is a design at the muonic exhaust velocity delivering exactly $T$ ---
thrust is bought with mass flow and power.  Consequently the engineering
question is a question about power budgets, and is recorded as a proposition

\begin{verbatim}
ThrustOptimizationConjecture T P  :=
  exists a design with exhaust velocity the muonic ceiling,
  thrust T, and jet power at most P
\end{verbatim}

\noindent which the development never asserts.  It is shown to be non-vacuous
(\texttt{thrust\_optimization\_satisfiable}) and to be settled by a power budget
rather than by the geometry (\texttt{thrust\_optimization\_is\_empirical}).  That
is the precise sense in which ``empirical thrust optimisations remain
unproven'': the framework's contribution is the hyperbola $T = 2P/v_e$, and
everything beyond it is engineering.


\section{Verification}
\label{sec:verification}

\noindent\textbf{Overall:} \textcolor{okgreen}{\textbf{pass} --- sources elaborate, every constant agrees, no forbidden token, no non-standard axiom}

\paragraph{1. Hygiene.}
The sources are scanned for \texttt{sorry}, \texttt{admit}, \texttt{axiom}
and \texttt{native\_decide} before anything is compiled; a hit aborts the
build.  Hits: \textbf{none}.

\paragraph{2. Elaboration.}
Each module is elaborated with \texttt{lean} in dependency order.
Total elaboration time: 13.06 s.

\begin{center}\small
\begin{tabular}{@{}lccc@{}}
\toprule
Module & Status & Time & Result \\
\midrule
\texttt{Basic} & elaborated & 2.13 s & ok \\
\texttt{Leptons} & elaborated & 2.12 s & ok \\
\texttt{Dimensions} & elaborated & 2.08 s & ok \\
\texttt{Velocity} & elaborated & 2.43 s & ok \\
\texttt{Regimes} & elaborated & 2.42 s & ok \\
\texttt{TheoreticalFramework} & elaborated & 1.88 s & ok \\
\bottomrule
\end{tabular}
\end{center}

\paragraph{3. Axiom audit.}
\texttt{\#print axioms} is run on every theorem and definition.  Grouped by
the set of axioms each declaration depends on:

\begin{center}\small
\begin{tabular}{@{}lc@{}}
\toprule
Axioms & Declarations \\
\midrule
none & 6 \\
\texttt{propext}, \texttt{Classical.choice}, \texttt{Quot.sound} & 133 \\
\bottomrule
\end{tabular}
\end{center}

\noindent None of these is \texttt{sorryAx}, and no declaration carries a
framework-specific axiom: the \texttt{propext} / \texttt{Classical.choice} /
\texttt{Quot.sound} triple is Lean's standard classical base, used by the
decision procedures, not by the mathematics of the framework.

\paragraph{4. Independent recomputation.}
Every numerical claim is recomputed in Python from its definition --- exact
integers and \texttt{Fraction} where the quantity is rational --- and compared
with the interval or numeral that Lean actually states.

{\footnotesize
\begin{longtable}{@{}p{0.325\textwidth}p{0.245\textwidth}p{0.175\textwidth}p{0.155\textwidth}p{0.05\textwidth}@{}}
\toprule
Theorem & Quantity & Lean & Recomputed & \\
\midrule
\endhead
\texttt{\scriptsize collisionVelocityMuonMax\_\allowbreak{}bounds} & muonic collision ceiling v₂, m/s & $[6{,}944,\ 6{,}945]$ & $6{,}944$ & \textcolor{okgreen}{ok} \\
\texttt{\scriptsize casimir\_\allowbreak{}amplification\_\allowbreak{}muon\_\allowbreak{}order} & casimir amplification (m\_\allowbreak{}μ/m\_\allowbreak{}e)⁵ & $[3.8 \times 10^{11},\ 4 \times 10^{11}]$ & $3.800596 \times 10^{11}$ & \textcolor{okgreen}{ok} \\
\texttt{\scriptsize casimir\_\allowbreak{}amplification\_\allowbreak{}tauon\_\allowbreak{}order} & casimir amplification (m\_\allowbreak{}τ/m\_\allowbreak{}e)⁵ & $[5.08 \times 10^{17},\ 5.1 \times 10^{17}]$ & $5.081869 \times 10^{17}$ & \textcolor{okgreen}{ok} \\
\texttt{\scriptsize casimir\_\allowbreak{}amplification\_\allowbreak{}tauon\_\allowbreak{}over\_\allowbreak{}muon} & (m\_\allowbreak{}τ/m\_\allowbreak{}μ)⁵ & $[1.3 \times 10^{6},\ 1.34 \times 10^{6}]$ & $1.337124 \times 10^{6}$ & \textcolor{okgreen}{ok} \\
\texttt{\scriptsize muonAmplification\_\allowbreak{}order} & velocity amplification 207³ & $[8.8 \times 10^{6},\ 8.9 \times 10^{6}]$ & $8.869743 \times 10^{6}$ & \textcolor{okgreen}{ok} \\
\texttt{\scriptsize tauonAmplification\_\allowbreak{}order} & velocity amplification 3477³ & $[4.2 \times 10^{10},\ 4.21 \times 10^{10}]$ & $4.203529 \times 10^{10}$ & \textcolor{okgreen}{ok} \\
\texttt{\scriptsize tauon\_\allowbreak{}over\_\allowbreak{}muon\_\allowbreak{}amplification} & 3477³/207³ & $[4{,}730,\ 4{,}740]$ & $4{,}739$ & \textcolor{okgreen}{ok} \\
\texttt{\scriptsize muonMaxEjection\_\allowbreak{}bounds} & muonic ejection ceiling, m/s & $[2.0682 \times 10^{7},\ 2.06821 \times 10^{7}]$ & $2.068202 \times 10^{7}$ & \textcolor{okgreen}{ok} \\
\texttt{\scriptsize muonMaxEjection\_\allowbreak{}km\_\allowbreak{}per\_\allowbreak{}s} & muonic ejection ceiling, km/s & $[20{,}680,\ 20{,}700]$ & $20{,}682$ & \textcolor{okgreen}{ok} \\
\texttt{\scriptsize muonSpecificImpulse\_\allowbreak{}bounds} & muonic specific impulse, s & $[2.108 \times 10^{6},\ 2.109 \times 10^{6}]$ & $2.108979 \times 10^{6}$ & \textcolor{okgreen}{ok} \\
\texttt{\scriptsize muonImpulse\_\allowbreak{}over\_\allowbreak{}raptor} & Isp(muonic)/Isp(Raptor) & $[5{,}800,\ 5{,}820]$ & $5{,}810$ & \textcolor{okgreen}{ok} \\
\texttt{\scriptsize tauonMaxEjection\_\allowbreak{}bounds} & tauonic ejection ceiling, m/s & $[4.2403 \times 10^{12},\ 4.2404 \times 10^{12}]$ & $4.240334 \times 10^{12}$ & \textcolor{okgreen}{ok} \\
\texttt{\scriptsize tauonEntryEjection\_\allowbreak{}bounds} & tauonic ejection at entry, m/s & $[1.4237 \times 10^{9},\ 1.4239 \times 10^{9}]$ & $1.423785 \times 10^{9}$ & \textcolor{okgreen}{ok} \\
\texttt{\scriptsize tauonMedianEjection\_\allowbreak{}bounds} & tauonic ejection at median input, m/s & $[2.8475 \times 10^{10},\ 2.8476 \times 10^{10}]$ & $2.847571 \times 10^{10}$ & \textcolor{okgreen}{ok} \\
\texttt{\scriptsize source\_\allowbreak{}tauon\_\allowbreak{}max\_\allowbreak{}velocity\_\allowbreak{}underestimates} & (verified/source)² at the tauonic ceiling & $[10.01,\ 10.02]$ & $10.0136$ & \textcolor{okgreen}{ok} \\
\texttt{\scriptsize source\_\allowbreak{}tauon\_\allowbreak{}max\_\allowbreak{}velocity\_\allowbreak{}factor} & verified/source at the tauonic ceiling & $[3.164,\ 3.165]$ & $3.1644$ & \textcolor{okgreen}{ok} \\
\texttt{\scriptsize muonic\_\allowbreak{}thrust\_\allowbreak{}to\_\allowbreak{}power\_\allowbreak{}bounds} & muonic thrust per watt, N/W & $[9.6 \times 10^{-8},\ 9.7 \times 10^{-8}]$ & $9.670234 \times 10^{-8}$ & \textcolor{okgreen}{ok} \\
\texttt{\scriptsize muonic\_\allowbreak{}thrust\_\allowbreak{}to\_\allowbreak{}power\_\allowbreak{}shortfall} & v\_\allowbreak{}e(muonic)/v\_\allowbreak{}e(Raptor) & $[5{,}800,\ 5{,}820]$ & $5{,}810$ & \textcolor{okgreen}{ok} \\
\texttt{\scriptsize domainBoundary\_\allowbreak{}km\_\allowbreak{}per\_\allowbreak{}s} & domain boundary, km/s & $[20{,}680,\ 20{,}700]$ & $20{,}682$ & \textcolor{okgreen}{ok} \\
\texttt{\scriptsize massRatioMuon\_\allowbreak{}pow\_\allowbreak{}five} & 207⁵ & exact numeral & \texttt{380059617807} & \textcolor{okgreen}{ok} \\
\texttt{\scriptsize massRatioTauon\_\allowbreak{}pow\_\allowbreak{}five} & 3477⁵ & exact numeral & \texttt{508186885683281157} & \textcolor{okgreen}{ok} \\
\texttt{\scriptsize muonAmplification\_\allowbreak{}eq} & 207³ & exact numeral & \texttt{8869743} & \textcolor{okgreen}{ok} \\
\texttt{\scriptsize tauonAmplification\_\allowbreak{}eq} & 3477³ & exact numeral & \texttt{42035292333} & \textcolor{okgreen}{ok} \\
\texttt{\scriptsize collisionVelocityMuonMax\_\allowbreak{}eq} & 25 000 km/h in m/s & exact numeral & \texttt{62500 / 9} & \textcolor{okgreen}{ok} \\
\bottomrule
\end{longtable}
}

\paragraph{5. Provenance.}
Source digests (SHA-256, first 16 hex digits) at report time:

\begin{center}\small
\begin{tabular}{@{}lc@{}}
\toprule
Module & Digest \\
\midrule
\texttt{Basic} & \texttt{b15febfb52f101a4} \\
\texttt{Leptons} & \texttt{9ac624893fc8a88b} \\
\texttt{Dimensions} & \texttt{490237f0691b40ab} \\
\texttt{Velocity} & \texttt{d1716d3817282452} \\
\texttt{Regimes} & \texttt{0ce9852307164c64} \\
\texttt{TheoreticalFramework} & \texttt{bc1e1b637f57f99e} \\
\bottomrule
\end{tabular}
\end{center}

\noindent Report generated 2026-11-05 13:39 UTC.


\clearpage
\appendix
\section{Declaration index}
\label{app:index}

Every declaration of the development, with its docstring.  The list is produced
by parsing the sources, so no declaration can be omitted by an oversight: if it
compiles, it is here.

\begin{center}\footnotesize
\begin{longtable}{@{}p{0.36\textwidth}p{0.10\textwidth}p{0.44\textwidth}@{}}
\toprule
Name & Kind & Content \\
\midrule
\endhead
\texttt{le\_\allowbreak{}of\_\allowbreak{}sq\_\allowbreak{}le\_\allowbreak{}sq} & \texttt{theorem} & \small Comparison of two non-negative reals through their squares.  This is the tool that turns an exactly computed \texttt{vf²} into a bound on \texttt{vf}. \\
\texttt{effRadius} & \texttt{def} & \small Effective radius of a lepton, \texttt{r\_\allowbreak{}ℓ = κ / m\_\allowbreak{}ℓ}. \texttt{κ} has dimensions of mass × length.  The framework fixes \texttt{κ} by requiring that the electron's effective radius reproduces its classical radius, \texttt{r\_\allowbreak{}e = 2.8 × 10⁻¹⁵ m}; no numerical value is needed for the scaling laws, which are \texttt{κ}-free. \\
\texttt{effRadius\_\allowbreak{}pos} & \texttt{theorem} & \small The effective radius of a positive-mass lepton is positive. \\
\texttt{effRadius\_\allowbreak{}div} & \texttt{theorem} & \small Effective radii are inversely proportional to mass: the \emph{compaction} property that drives the whole framework. \\
\texttt{effRadius\_\allowbreak{}strictAnti} & \texttt{theorem} & \small A heavier lepton is the more compact one. \\
\texttt{casimirEnergy} & \texttt{def} & \small Casimir energy of the internal two-plate geometry, in units where the prefactor \texttt{π²ℏc/720} is absorbed into the area coefficient. \texttt{E\_\allowbreak{}C $\propto$ --- S / d³} with \texttt{S} the plate area and \texttt{d = r\_\allowbreak{}ℓ = κ/m} the plate separation.  The parallel-plate area ansatz \texttt{S $\propto$ m²} is the source's, and is recorded here explicitly by the \texttt{m\textasciicircum{}2} numerator. \\
\texttt{casimirEnergy\_\allowbreak{}eq\_\allowbreak{}quintic} & \texttt{theorem} & \small \textbf{Quintic law.}  Under the inverse-radius ansatz the Casimir energy of a lepton is proportional to the fifth power of its mass. \\
\texttt{casimirEnergy\_\allowbreak{}ratio} & \texttt{theorem} & \small \textbf{Amplification ratio.}  Between two leptons, the Casimir amplification is the fifth power of their mass ratio.  The constant \texttt{κ} cancels. \\
\texttt{casimirEnergy\_\allowbreak{}ratio\_\allowbreak{}radii} & \texttt{theorem} & \small The same amplification, read geometrically: it is the fifth power of the \emph{inverse} ratio of effective radii.  Compactness and vacuum coupling are two faces of one quantity. \\
\texttt{CasimirAnsatz} & \texttt{structure} & \small The geometric ansatz of the framework: an internal plate geometry whose separation is the effective radius and whose area scales as the square of the mass. \\
\texttt{CasimirAnsatz.plateArea} & \texttt{def} & \small Plate area of the internal geometry, \texttt{S = areaCoeff · m²}. \\
\texttt{CasimirAnsatz.energy} & \texttt{def} & \small Casimir energy of the ansatz, \texttt{E = -S/d³} up to a positive prefactor; its magnitude is \texttt{S/d³}. \\
\texttt{CasimirAnsatz.signedEnergy} & \texttt{def} & \small The signed (physical, attractive) Casimir energy, \texttt{E = -C · S/d³} with \texttt{C \textgreater{} 0}. \\
\texttt{CasimirAnsatz.energy\_\allowbreak{}eq\_\allowbreak{}quintic} & \texttt{theorem} & \small  \\
\texttt{CasimirAnsatz.energy\_\allowbreak{}ratio} & \texttt{theorem} & \small The amplification ratio of the ansatz is the fifth power of the mass ratio. \\
\texttt{CasimirAnsatz.energy\_\allowbreak{}ratio\_\allowbreak{}parameter\_\allowbreak{}free} & \texttt{theorem} & \small \textbf{Parameter-freedom of the prediction.}  Two arbitrary choices of the framework's free constants — \texttt{κ}, the area coefficient, even the overall Casimir prefactor — give the same amplification ratio.  The ratio \texttt{(m₁/m₂)⁵} is thus a genuine prediction of the ansatz rather than a fit. \\
\texttt{CasimirAnsatz.signedEnergy\_\allowbreak{}ratio} & \texttt{theorem} & \small The sign convention of the Casimir effect does not affect the amplification: the physical energy is negative (attractive), and the ratio is unchanged. \\
\texttt{massElectronMeV} & \texttt{def} & \small Electron mass, MeV/c². \\
\texttt{massMuonMeV} & \texttt{def} & \small Muon mass, MeV/c². \\
\texttt{massTauonMeV} & \texttt{def} & \small Tauon mass, MeV/c². \\
\texttt{massRatioMuon} & \texttt{def} & \small Muon-to-electron mass ratio, as used throughout the framework. \\
\texttt{massRatioTauon} & \texttt{def} & \small Tauon-to-electron mass ratio, as used throughout the framework. \\
\texttt{massRatioMuon\_\allowbreak{}pos} & \texttt{theorem} & \small  \\
\texttt{massRatioTauon\_\allowbreak{}pos} & \texttt{theorem} & \small  \\
\texttt{massRatioMuon\_\allowbreak{}consistent} & \texttt{theorem} & \small The integer ratio \texttt{207} reproduces the tabulated muon/electron mass ratio to better than a quarter of a unit. \\
\texttt{massRatioTauon\_\allowbreak{}consistent} & \texttt{theorem} & \small The integer ratio \texttt{3477} reproduces the tabulated tauon/electron mass ratio to better than half a unit. \\
\texttt{casimir\_\allowbreak{}amplification\_\allowbreak{}muon\_\allowbreak{}general} & \texttt{theorem} & \small \textbf{Muonic Casimir amplification}, in the form \texttt{E\_\allowbreak{}μ / E\_\allowbreak{}e = (m\_\allowbreak{}μ/m\_\allowbreak{}e)⁵}. \\
\texttt{casimir\_\allowbreak{}amplification\_\allowbreak{}tauon\_\allowbreak{}general} & \texttt{theorem} & \small \textbf{Tauonic Casimir amplification}, in the form \texttt{E\_\allowbreak{}τ / E\_\allowbreak{}e = (m\_\allowbreak{}τ/m\_\allowbreak{}e)⁵}. \\
\texttt{massRatioMuon\_\allowbreak{}pow\_\allowbreak{}five} & \texttt{theorem} & \small Exact value of the muonic amplification, \texttt{207⁵ = 380 059 617 807}. \\
\texttt{massRatioTauon\_\allowbreak{}pow\_\allowbreak{}five} & \texttt{theorem} & \small Exact value of the tauonic amplification, \texttt{3477⁵ = 508 186 885 683 281 157}. \\
\texttt{casimir\_\allowbreak{}amplification\_\allowbreak{}muon\_\allowbreak{}order} & \texttt{theorem} & \small \texttt{(207)⁵ ≈ 4 × 10¹¹}: the amplification lies between \texttt{3.8 × 10¹¹} and \texttt{4 × 10¹¹}. \\
\texttt{casimir\_\allowbreak{}amplification\_\allowbreak{}tauon\_\allowbreak{}order} & \texttt{theorem} & \small \texttt{(3477)⁵ ≈ 5.1 × 10¹⁷}: the amplification lies between \texttt{5.08 × 10¹⁷} and \texttt{5.1 × 10¹⁷}. \\
\texttt{casimir\_\allowbreak{}amplification\_\allowbreak{}tauon\_\allowbreak{}over\_\allowbreak{}muon} & \texttt{theorem} & \small The tauon amplifies the vacuum coupling by a further factor \texttt{(3477/207)⁵ ≈ 1.3 × 10⁶} over the muon. \\
\texttt{casimir\_\allowbreak{}amplification\_\allowbreak{}hierarchy} & \texttt{theorem} & \small The amplification is monotone in the mass ratio: the lepton hierarchy \texttt{e \textless{} μ \textless{} τ} is a strict hierarchy of vacuum couplings. \\
\texttt{fieldAmplificationSource} & \texttt{def} & \small Equivalent-field amplification under the source's volume ansatz \texttt{V $\propto$ S·d³}: \texttt{B $\propto$ m³}. \\
\texttt{fieldAmplificationPlate} & \texttt{def} & \small Equivalent-field amplification under the literal plate volume \texttt{V = S·d}: \texttt{B $\propto$ m²}. \\
\texttt{field\_\allowbreak{}amplification\_\allowbreak{}source} & \texttt{theorem} & \small Source ansatz: \texttt{B² $\propto$ E\_\allowbreak{}C / (S·d³) $\propto$ m⁶}, hence \texttt{B $\propto$ m³}. \\
\texttt{field\_\allowbreak{}amplification\_\allowbreak{}plate} & \texttt{theorem} & \small Literal plate volume: \texttt{B² $\propto$ E\_\allowbreak{}C / (S·d) $\propto$ m⁴}, hence \texttt{B $\propto$ m²}. \\
\texttt{field\_\allowbreak{}amplification\_\allowbreak{}muon\_\allowbreak{}order} & \texttt{theorem} & \small The source's field amplification for the muon, \texttt{(207)³ ≈ 9 × 10⁶}. \\
\texttt{casimirMassExponent} & \texttt{def} & \small Exponent of the mass in the Casimir energy of the compactness ansatz. \\
\texttt{generalizedDimension} & \texttt{def} & \small The generalized dimension, fixed by the quintic mass dependence: \texttt{n = 5}. \\
\texttt{generalizedDimension\_\allowbreak{}eq\_\allowbreak{}five} & \texttt{theorem} & \small The generalized dimension is five. \\
\texttt{structuralTerm} & \texttt{def} & \small The structural part of the anomalous moment, \texttt{a\_\allowbreak{}str = C · m\textasciicircum{}n}. \\
\texttt{anomalousMoment} & \texttt{def} & \small The framework's decomposition \texttt{a\_\allowbreak{}ℓ = a\_\allowbreak{}QED + a\_\allowbreak{}EW+had + a\_\allowbreak{}str}. \\
\texttt{structuralTerm\_\allowbreak{}quintic} & \texttt{theorem} & \small The structural term is quintic in the mass. \\
\texttt{structuralTerm\_\allowbreak{}ratio} & \texttt{theorem} & \small The structural term amplifies between two leptons by the fifth power of the mass ratio: the same number as the Casimir amplification, which is the reason a single exponent serves both. \\
\texttt{velocityExponent} & \texttt{def} & \small Exponent of the mass ratio in the velocity law, \texttt{n − 2}. \\
\texttt{velocityExponent\_\allowbreak{}eq\_\allowbreak{}three} & \texttt{theorem} & \small With \texttt{n = 5} the velocity law carries the cube of the mass ratio. \\
\texttt{velocityExponent\_\allowbreak{}eq} & \texttt{theorem} & \small The velocity exponent is \texttt{n − 2} for the framework's \texttt{n}. \\
\texttt{generalised\_\allowbreak{}unit\_\allowbreak{}exponents} & \texttt{theorem} & \small Unit bookkeeping of \texttt{[vₙ] = m\textasciicircum{}(2/n)/s}: the mass exponent balances to \texttt{2}. \\
\texttt{time\_\allowbreak{}exponent\_\allowbreak{}balance} & \texttt{theorem} & \small Unit bookkeeping of the time factor \texttt{dt₀\textasciicircum{}(n−2)}: the exponents close on \texttt{n}. \\
\texttt{GeneralisedVelocityRelation} & \texttt{def} & \small The generalized-velocity relation \texttt{vₙⁿ · dt₀\textasciicircum{}(n−2) = v₂²} of the framework. \\
\texttt{scale\_\allowbreak{}substitution} & \texttt{theorem} & \small \textbf{Scale substitution.}  Replacing the time scale by \texttt{q·dt₀} multiplies the relation by \texttt{q\textasciicircum{}(n−2)}, which is the framework's substitution rule \texttt{vₙ·(q·dt₀)\textasciicircum{}(n−2) = v₂²·q\textasciicircum{}(n−2)}; the parameter \texttt{q = m\_\allowbreak{}ℓ/m₀} is what links the muonic and tauonic regimes. \\
\texttt{generalisedVelocity} & \texttt{def} & \small Ejection velocity of particles produced in lepton–lepton collisions in the generalized framework: the non-negative root of \texttt{vf² = v₂²·r\textasciicircum{}(n−2)}, where \texttt{r} is the mass ratio \texttt{m\_\allowbreak{}ℓ/m₀}. \\
\texttt{velocity\_\allowbreak{}law} & \texttt{theorem} & \small \textbf{The velocity law}, for arbitrary generalized dimension \texttt{n}: \texttt{vf² = v₂² · (m\_\allowbreak{}ℓ/m₀)\textasciicircum{}(n−2)}. \\
\texttt{generalisedVelocity\_\allowbreak{}eq\_\allowbreak{}mul} & \texttt{theorem} & \small The paper's form of the law, \texttt{vf = v₂ · √((m\_\allowbreak{}ℓ/m₀)\textasciicircum{}(n−2))}, agrees with the definition for non-negative input velocity. \\
\texttt{generalisedVelocity\_\allowbreak{}nonneg} & \texttt{theorem} & \small Ejection velocities are non-negative. \\
\texttt{generalisedVelocity\_\allowbreak{}bounds} & \texttt{theorem} & \small Two-sided bounds on an ejection velocity, from exact bounds on \texttt{vf²}. This is the workhorse that turns a computed square into a velocity interval without any rounding step. \\
\texttt{velocity\_\allowbreak{}law\_\allowbreak{}five} & \texttt{theorem} & \small The velocity law at the framework's generalized dimension: \texttt{vf² = v₂² · (m\_\allowbreak{}ℓ/m₀)³}. \\
\texttt{generalisedVelocity\_\allowbreak{}five\_\allowbreak{}eq\_\allowbreak{}mul} & \texttt{theorem} & \small The law's explicit linear-in-\texttt{v₂} form at \texttt{n = 5}: \texttt{vf = v₂·√(r³)}. \\
\texttt{generalisedVelocity\_\allowbreak{}mono} & \texttt{theorem} & \small The ejection velocity grows with the lepton mass ratio: heavier collision partners eject faster. \\
\texttt{generalisedVelocity\_\allowbreak{}mono\_\allowbreak{}v₂} & \texttt{theorem} & \small The ejection velocity grows with the collision velocity. \\
\texttt{generalisedVelocity\_\allowbreak{}zero} & \texttt{theorem} & \small The ejection velocity vanishes with the collision velocity: no collisions, no propulsion. \\
\texttt{kmhToMs} & \texttt{def} & \small Kilometres per hour to metres per second, \texttt{1 km/h = 5/18 m/s}. \\
\texttt{kmhToMs\_\allowbreak{}eq} & \texttt{theorem} & \small The conversion agrees with the definitional \texttt{1000/3600}. \\
\texttt{standardGravity} & \texttt{def} & \small Standard gravity, \texttt{g₀ = 9.80665 m/s²}, used to express specific impulse. \\
\texttt{standardGravity\_\allowbreak{}pos} & \texttt{theorem} & \small  \\
\texttt{speedOfLight} & \texttt{def} & \small Speed of light in vacuum, m/s. \\
\texttt{specificImpulse} & \texttt{def} & \small Specific impulse, \texttt{Isp = v\_\allowbreak{}e / g₀}, in seconds. \\
\texttt{raptorSpecificImpulse} & \texttt{def} & \small Vacuum specific impulse of a Raptor engine, seconds: the chemical benchmark of the source text. \\
\texttt{collisionVelocityMuonMax} & \texttt{def} & \small Upper end of the muonic input window, \texttt{v₂ = 25 000 km/h}. \\
\texttt{collisionVelocityMuonMax\_\allowbreak{}eq} & \texttt{theorem} & \small The muonic input ceiling in m/s. \\
\texttt{collisionVelocityMuonMax\_\allowbreak{}bounds} & \texttt{theorem} & \small \texttt{25 000 km/h ≈ 6 944 m/s}. \\
\texttt{muonAmplification} & \texttt{def} & \small Muonic amplification factor, \texttt{(m\_\allowbreak{}μ/m\_\allowbreak{}e)\textasciicircum{}(n−2)}. \\
\texttt{tauonAmplification} & \texttt{def} & \small Tauonic amplification factor, \texttt{(m\_\allowbreak{}τ/m\_\allowbreak{}e)\textasciicircum{}(n−2)}. \\
\texttt{muonAmplification\_\allowbreak{}eq} & \texttt{theorem} & \small \texttt{207³ = 8 869 743}. \\
\texttt{tauonAmplification\_\allowbreak{}eq} & \texttt{theorem} & \small \texttt{3477³ = 42 035 292 333}. \\
\texttt{muonAmplification\_\allowbreak{}order} & \texttt{theorem} & \small The muonic amplification is \texttt{(207)³ ≈ 8.9 × 10⁶}. \\
\texttt{tauonAmplification\_\allowbreak{}order} & \texttt{theorem} & \small The tauonic amplification is \texttt{(3477)³ ≈ 4.2 × 10¹⁰}. \\
\texttt{tauon\_\allowbreak{}over\_\allowbreak{}muon\_\allowbreak{}amplification} & \texttt{theorem} & \small The tauonic amplification exceeds the muonic one by \texttt{(3477/207)³ ≈ 4.74 × 10³}. \\
\texttt{muonMaxEjection} & \texttt{def} & \small Maximum ejection velocity of the muonic regime, at \texttt{v₂ = 25 000 km/h}. \\
\texttt{muonMaxEjection\_\allowbreak{}sq} & \texttt{theorem} & \small The ejection velocity obeys the law: \texttt{vf² = v₂² · 207³}. \\
\texttt{muonMaxEjection\_\allowbreak{}bounds} & \texttt{theorem} & \small \texttt{vf ≈ 2.06820 × 10⁷ m/s} for the muonic ceiling. \\
\texttt{muonMaxEjection\_\allowbreak{}km\_\allowbreak{}per\_\allowbreak{}s} & \texttt{theorem} & \small The source's \texttt{≈ 20 700 km/s}, proved as an interval in km/s. \\
\texttt{muonMaxEjection\_\allowbreak{}subluminal} & \texttt{theorem} & \small The muonic ejection velocity is subluminal: the muonic regime is compatible with ordinary special relativity. \\
\texttt{collisionVelocityMuonMax\_\allowbreak{}subluminal} & \texttt{theorem} & \small The muonic input window is subluminal as well. \\
\texttt{muonSpecificImpulse\_\allowbreak{}bounds} & \texttt{theorem} & \small Specific impulse of the muonic ceiling, \texttt{Isp = v\_\allowbreak{}e/g₀ ≈ 2.109 × 10⁶ s}. \\
\texttt{muonImpulse\_\allowbreak{}over\_\allowbreak{}raptor} & \texttt{theorem} & \small The muonic impulse exceeds the Raptor benchmark by a factor of about \texttt{5.8 × 10³}, i.e. "roughly 6 000 times" as the source states. \\
\texttt{tauonicInputMax} & \texttt{def} & \small Input ceiling of the tauonic regime: the muonic regime's maximum ejection velocity, \texttt{≈ 20 682 km/s}.  Taking the unrounded value is what makes the two regimes exactly continuous. \\
\texttt{tauonMaxEjection} & \texttt{def} & \small Maximum ejection velocity of the tauonic regime, at \texttt{v₂ = tauonicInputMax}. \\
\texttt{tauonMaxEjection\_\allowbreak{}sq} & \texttt{theorem} & \small The tauonic ejection velocity obeys the law: \texttt{vf² = v₂² · 3477³}. \\
\texttt{tauonMaxEjection\_\allowbreak{}bounds} & \texttt{theorem} & \small \texttt{vf ≈ 4.2403 × 10¹² m/s} for the tauonic ceiling. \\
\texttt{tauonMaxEjection\_\allowbreak{}superluminal} & \texttt{theorem} & \small The tauonic ceiling is superluminal: the framework's rejection of the Lorentz factor is not optional decoration, it is what its own numbers require. \\
\texttt{tauonicInputMax\_\allowbreak{}subluminal} & \texttt{theorem} & \small The tauonic input window is subluminal; only the output exceeds \texttt{c}. \\
\texttt{sourceTauonMaxVelocity} & \texttt{def} & \small The tauonic ceiling ejection velocity as printed in the source text. \\
\texttt{sourceTauonMedianVelocity} & \texttt{def} & \small The tauonic median ejection velocity as printed in the source text. \\
\texttt{source\_\allowbreak{}tauon\_\allowbreak{}max\_\allowbreak{}velocity\_\allowbreak{}underestimates} & \texttt{theorem} & \small The source figure, squared, falls short of the law by a factor of \texttt{10.01}: the printed value is \texttt{√10} too small. \\
\texttt{source\_\allowbreak{}tauon\_\allowbreak{}max\_\allowbreak{}velocity\_\allowbreak{}factor} & \texttt{theorem} & \small Consequently the source figure underestimates the ejection velocity itself by the factor \texttt{√10 ≈ 3.164}. \\
\texttt{tauonEntryEjection} & \texttt{def} & \small Tauonic ejection velocity at the entry velocity of the tauonic window. \\
\texttt{tauonEntryEjection\_\allowbreak{}bounds} & \texttt{theorem} & \small \texttt{vf ≈ 1.4238 × 10⁹ m/s} at the tauonic entry velocity: the source's \texttt{1.42 × 10⁹} is correct here, because this \texttt{vf²} happens to carry an even decimal exponent. \\
\texttt{tauonMedianEjection} & \texttt{def} & \small Tauonic ejection velocity at a median input of \texttt{500 000 km/h}. \\
\texttt{tauonMedianEjection\_\allowbreak{}bounds} & \texttt{theorem} & \small \texttt{vf ≈ 2.8477 × 10¹⁰ m/s} at the median input — not the \texttt{9.0 × 10¹⁰} printed in the source, which carries the same \texttt{√10} slip. \\
\texttt{amplification\_\allowbreak{}mul} & \texttt{theorem} & \small The amplification factor is multiplicative in the mass ratio: chaining two mass scales multiplies the amplification.  This is what makes the regimes composable rather than merely analogous. \\
\texttt{amplification\_\allowbreak{}strictMono} & \texttt{theorem} & \small The amplification is strictly increasing in the mass ratio on the positive axis: heavier collision partners amplify more. \\
\texttt{tauon\_\allowbreak{}amplification\_\allowbreak{}factors} & \texttt{theorem} & \small The tauonic amplification is the muonic one times the amplification of the residual ratio \texttt{m\_\allowbreak{}τ/m\_\allowbreak{}μ}: \texttt{3477³ = 207³ · (3477/207)³}. \\
\texttt{generalisedVelocity\_\allowbreak{}ge\_\allowbreak{}input} & \texttt{theorem} & \small For a mass ratio at least one, the ejection velocity is at least the collision velocity.  This is what lets the tauonic window start where the muonic window ends. \\
\texttt{one\_\allowbreak{}le\_\allowbreak{}massRatioMuon} & \texttt{theorem} & \small  \\
\texttt{one\_\allowbreak{}le\_\allowbreak{}massRatioTauon} & \texttt{theorem} & \small  \\
\texttt{PropulsiveRegime} & \texttt{structure} & \small A propulsive regime of the framework: a charged-lepton mass ratio, the window of admissible collision velocities, and the law-induced quantities derived from them. \\
\texttt{PropulsiveRegime.amplification} & \texttt{def} & \small Amplification factor of the regime, \texttt{(m\_\allowbreak{}ℓ/m\_\allowbreak{}e)\textasciicircum{}(n−2)}. \\
\texttt{PropulsiveRegime.maxOutput} & \texttt{def} & \small Maximum ejection velocity of the regime, at the top of its input window. \\
\texttt{PropulsiveRegime.window} & \texttt{def} & \small The input window as a set. \\
\texttt{PropulsiveRegime.maxOutput\_\allowbreak{}sq} & \texttt{theorem} & \small The regime's maximum output obeys the velocity law. \\
\texttt{PropulsiveRegime.inputUpper\_\allowbreak{}le\_\allowbreak{}maxOutput} & \texttt{theorem} & \small The maximum output is at least the input ceiling. \\
\texttt{PropulsiveRegime.continuous\_\allowbreak{}maxOutput} & \texttt{theorem} & \small Within a regime the velocity law is continuous in the collision velocity. \\
\texttt{PropulsiveRegime.continuous\_\allowbreak{}in\_\allowbreak{}massRatio} & \texttt{theorem} & \small Within a regime the velocity law is continuous in the mass ratio. \\
\texttt{muonicRegime} & \texttt{def} & \small \textbf{The muonic regime.}  Muon–muon collisions, mass ratio \texttt{207}, collision velocities up to \texttt{25 000 km/h}; replaces chemical propellants. \\
\texttt{tauonicRegime} & \texttt{def} & \small \textbf{The tauonic regime.}  Tauon–tauon collisions, mass ratio \texttt{3477}, input window from the muonic ceiling up to the muonic regime's output ceiling. \\
\texttt{muonicRegime\_\allowbreak{}amplification} & \texttt{theorem} & \small The muonic amplification, through the structure. \\
\texttt{tauonicRegime\_\allowbreak{}amplification} & \texttt{theorem} & \small The tauonic amplification, through the structure. \\
\texttt{regimes\_\allowbreak{}strictly\_\allowbreak{}ordered} & \texttt{theorem} & \small The regime hierarchy is strict: the tauonic amplification exceeds the muonic one by \texttt{≈ 4.74 × 10³}. \\
\texttt{muonicRegime\_\allowbreak{}maxOutput} & \texttt{theorem} & \small The muonic regime's output ceiling, through the structure. \\
\texttt{tauonicRegime\_\allowbreak{}maxOutput} & \texttt{theorem} & \small The tauonic regime's output ceiling, through the structure. \\
\texttt{windows\_\allowbreak{}adjacent} & \texttt{theorem} & \small \textbf{Adjacency.}  The tauonic window begins exactly where the muonic window ends: \texttt{25 000 km/h}, the muonic collision ceiling. \\
\texttt{chain\_\allowbreak{}closes} & \texttt{theorem} & \small \textbf{The chain closes.}  The upper end of the tauonic \emph{input} window is the muonic regime's maximum \emph{output}.  This is the boundary the framework calls \texttt{v₂ = 20 700 km/s}, and taking the unrounded value is what makes the match exact rather than approximate. \\
\texttt{regimes\_\allowbreak{}tile} & \texttt{theorem} & \small \textbf{Tiling.}  The two input windows cover one interval with no gap and no overlap. \\
\texttt{tauonic\_\allowbreak{}window\_\allowbreak{}nested\_\allowbreak{}in\_\allowbreak{}output} & \texttt{theorem} & \small \textbf{Nesting.}  The tauonic window lies inside the muonic regime's output range: the second regime's inputs are the first regime's reachable outputs. \\
\texttt{domainBoundary} & \texttt{def} & \small \textbf{The domain boundary.}  The boundary of the source's §8.2, \texttt{v₂ = 20 700 km/s}, is the tauonic input ceiling, and it is the muonic output ceiling unrounded. \\
\texttt{domainBoundary\_\allowbreak{}eq\_\allowbreak{}muonic\_\allowbreak{}output} & \texttt{theorem} & \small The boundary is the muonic regime's maximum output. \\
\texttt{domainBoundary\_\allowbreak{}km\_\allowbreak{}per\_\allowbreak{}s} & \texttt{theorem} & \small The boundary is the muonic ejection ceiling, hence \texttt{≈ 20 700 km/s}. \\
\texttt{domainBoundary\_\allowbreak{}subluminal} & \texttt{theorem} & \small The boundary itself is subluminal; it is the tauonic \emph{output} that exceeds \texttt{c}, not the input ceiling that defines the hand-off. \\
\texttt{regime\_\allowbreak{}continuity} & \texttt{theorem} & \small \textbf{Continuity of the two-regime structure}, collected in one statement: the windows are adjacent, the chain closes on the muonic output ceiling, they tile an interval, and the law is continuous in the collision velocity throughout. \\
\texttt{ThrusterDesign} & \texttt{structure} & \small A thruster design obeying the framework's momentum and energy balance. \\
\texttt{raptorExhaustVelocity} & \texttt{def} & \small Vacuum exhaust velocity of the Raptor benchmark. \\
\texttt{thrust\_\allowbreak{}to\_\allowbreak{}power} & \texttt{theorem} & \small \textbf{The thrust–power tradeoff.}  For any design, \texttt{T/P = 2/v\_\allowbreak{}e} exactly.  High specific impulse costs thrust per watt, without exception. \\
\texttt{thrust\_\allowbreak{}eq\_\allowbreak{}power} & \texttt{theorem} & \small Thrust in terms of jet power and exhaust velocity. \\
\texttt{muonic\_\allowbreak{}thrust\_\allowbreak{}to\_\allowbreak{}power\_\allowbreak{}bounds} & \texttt{theorem} & \small At the muonic ceiling, thrust per watt is \texttt{≈ 9.67 × 10⁻⁸ N/W}. \\
\texttt{muonic\_\allowbreak{}thrust\_\allowbreak{}to\_\allowbreak{}power\_\allowbreak{}shortfall} & \texttt{theorem} & \small The muonic exhaust velocity is \texttt{≈ 5.8 × 10³} times the Raptor exhaust velocity, so its thrust per watt is that many times \emph{worse}: the specific-impulse gain and the thrust-per-watt loss are the same number. \\
\texttt{thrust\_\allowbreak{}not\_\allowbreak{}determined} & \texttt{theorem} & \small \textbf{Thrust is not determined by the geometry.}  At the exhaust velocity the leptonic kinematics fix, any positive thrust is attainable by a suitable design: thrust is bought with mass flow and power, not with the velocity law. \\
\texttt{ThrustOptimizationConjecture} & \texttt{def} & \small The engineering question the framework leaves open: at the exhaust velocity the geometry fixes, is there a design reaching a thrust target \texttt{T} within a power budget \texttt{P}?  Nothing in the geometric development decides it — see \texttt{thrust\_\allowbreak{}not\_\allowbreak{}determined} — so it is recorded, not proved. \\
\texttt{thrust\_\allowbreak{}optimization\_\allowbreak{}satisfiable} & \texttt{theorem} & \small The conjecture is not vacuous: a large enough power budget always suffices, because the thrust–power relation is a hyperbola and nothing else constrains it. \\
\texttt{thrust\_\allowbreak{}optimization\_\allowbreak{}is\_\allowbreak{}empirical} & \texttt{theorem} & \small Consequently the thrust question is \emph{empirical}, not geometric: the framework fixes \texttt{v\_\allowbreak{}e} and hence \texttt{T/P}, and leaves \texttt{(T, P)} free along that hyperbola.  This is the formal statement of "empirical thrust optimisations remain unproven". \\
\bottomrule
\end{longtable}
\end{center}


\clearpage
\section{The verified sources}
\label{app:sources}

The complete development, as elaborated by Lean.  Digests are in
\S\ref{sec:verification}.


\subsection*{\texttt{TheoreticalFramework/Basic.lean}}
\vspace{2pt}
{\ttfamily\scriptsize\setlength{\parindent}{0pt}\setlength{\baselineskip}{7.6pt}\setlength{\parskip}{0pt}\raggedright\obeyspaces\obeylines
/-
\# TheoreticalFramework.Basic
\mbox{}
Core derivations of the leptonic-compactness framework.
\mbox{}
The model under study attributes part of the anomalous magnetic moment of a
massive lepton to an *internal geometric compactness*: a lepton of mass `m`
carries an effective radius
\mbox{}
  `r\_\allowbreak{}ℓ = κ / m`,
\mbox{}
with `κ` a fixed constant of the framework (dimensions of mass × length).  The
internal geometry is modelled, following the source text, as a pair of plates
whose separation is `d = r\_\allowbreak{}ℓ` and whose area scales with the square of the mass,
`S $\propto$ m²`.  The Casimir energy of such a cavity scales as
\mbox{}
  `E\_\allowbreak{}C $\propto$ - S / d³`,
\mbox{}
hence as `- m² / r\_\allowbreak{}ℓ³ = - m⁵ / κ³`.
\mbox{}
The mathematical content proved here is the scaling: the Casimir energy of a
lepton is quintic in its mass, and therefore the amplification ratio between two
leptons depends only on the ratio of their masses — not on `κ`, not on the area
coefficient, and not on the overall Casimir constant.  That parameter-freedom is
what lets the framework issue a numerical prediction, and it is the reason the
derivation can be checked exactly rather than up to an unknown prefactor.
\mbox{}
The geometric *ansatz* itself (`S $\propto$ m²`, `d = r\_\allowbreak{}ℓ`) is recorded as an explicit
model assumption in `CasimirAnsatz`, not asserted as a derived fact.
-/
import Mathlib
\mbox{}
namespace TheoreticalFramework
\mbox{}
noncomputable section
\mbox{}
/-- Comparison of two non-negative reals through their squares.  This is the
tool that turns an exactly computed `vf²` into a bound on `vf`. -/
theorem le\_\allowbreak{}of\_\allowbreak{}sq\_\allowbreak{}le\_\allowbreak{}sq \{a b : ℝ\} (ha : 0 ≤ a) (hb : 0 ≤ b) (h : a \textasciicircum{} 2 ≤ b \textasciicircum{} 2) : a ≤ b := by
  have habs := sq\_\allowbreak{}le\_\allowbreak{}sq.mp h
  rwa [abs\_\allowbreak{}of\_\allowbreak{}nonneg ha, abs\_\allowbreak{}of\_\allowbreak{}nonneg hb] at habs
\mbox{}
/-! \#\# Effective radius -/
\mbox{}
/-- Effective radius of a lepton, `r\_\allowbreak{}ℓ = κ / m\_\allowbreak{}ℓ`.
\mbox{}
`κ` has dimensions of mass × length.  The framework fixes `κ` by requiring that
the electron's effective radius reproduces its classical radius, `r\_\allowbreak{}e = 2.8 × 10⁻¹⁵ m`;
no numerical value is needed for the scaling laws, which are `κ`-free. -/
def effRadius (κ m : ℝ) : ℝ := κ / m
\mbox{}
/-- The effective radius of a positive-mass lepton is positive. -/
theorem effRadius\_\allowbreak{}pos \{κ m : ℝ\} (hκ : 0 \textless{} κ) (hm : 0 \textless{} m) : 0 \textless{} effRadius κ m :=
  div\_\allowbreak{}pos hκ hm
\mbox{}
/-- Effective radii are inversely proportional to mass: the *compaction* property
that drives the whole framework. -/
theorem effRadius\_\allowbreak{}div \{κ m₁ m₂ : ℝ\} (hκ : κ ≠ 0) (h₁ : m₁ ≠ 0) (h₂ : m₂ ≠ 0) :
    effRadius κ m₁ / effRadius κ m₂ = m₂ / m₁ := by
  unfold effRadius
  field\_\allowbreak{}simp
\mbox{}
/-- A heavier lepton is the more compact one. -/
theorem effRadius\_\allowbreak{}strictAnti \{κ : ℝ\} (hκ : 0 \textless{} κ) :
    StrictAntiOn (effRadius κ) (Set.Ioi 0) := by
  intro m₁ hm₁ m₂ hm₂ hlt
  simp only [Set.mem\_\allowbreak{}Ioi] at hm₁ hm₂
  exact div\_\allowbreak{}lt\_\allowbreak{}div\_\allowbreak{}of\_\allowbreak{}pos\_\allowbreak{}left hκ hm₁ hlt
\mbox{}
/-! \#\# Casimir energy of the internal geometry -/
\mbox{}
/-- Casimir energy of the internal two-plate geometry, in units where the
prefactor `π²ℏc/720` is absorbed into the area coefficient.
\mbox{}
`E\_\allowbreak{}C $\propto$ - S / d³` with `S` the plate area and `d = r\_\allowbreak{}ℓ = κ/m` the plate
separation.  The parallel-plate area ansatz `S $\propto$ m²` is the source's, and is
recorded here explicitly by the `m\textasciicircum{}2` numerator. -/
def casimirEnergy (κ m : ℝ) : ℝ := m \textasciicircum{} 2 / (effRadius κ m) \textasciicircum{} 3
\mbox{}
/-- **Quintic law.**  Under the inverse-radius ansatz the Casimir energy of a
lepton is proportional to the fifth power of its mass. -/
theorem casimirEnergy\_\allowbreak{}eq\_\allowbreak{}quintic \{κ m : ℝ\} (hκ : κ ≠ 0) (hm : m ≠ 0) :
    casimirEnergy κ m = m \textasciicircum{} 5 / κ \textasciicircum{} 3 := by
  unfold casimirEnergy effRadius
  field\_\allowbreak{}simp
\mbox{}
/-- **Amplification ratio.**  Between two leptons, the Casimir amplification is
the fifth power of their mass ratio.  The constant `κ` cancels. -/
theorem casimirEnergy\_\allowbreak{}ratio \{κ : ℝ\} (hκ : κ ≠ 0) \{m₁ m₂ : ℝ\} (h₁ : m₁ ≠ 0) (h₂ : m₂ ≠ 0) :
    casimirEnergy κ m₁ / casimirEnergy κ m₂ = (m₁ / m₂) \textasciicircum{} 5 := by
  rw [casimirEnergy\_\allowbreak{}eq\_\allowbreak{}quintic hκ h₁, casimirEnergy\_\allowbreak{}eq\_\allowbreak{}quintic hκ h₂]
  field\_\allowbreak{}simp
\mbox{}
/-- The same amplification, read geometrically: it is the fifth power of the
*inverse* ratio of effective radii.  Compactness and vacuum coupling are two
faces of one quantity. -/
theorem casimirEnergy\_\allowbreak{}ratio\_\allowbreak{}radii \{κ m₁ m₂ : ℝ\} (hκ : κ ≠ 0) (h₁ : m₁ ≠ 0) (h₂ : m₂ ≠ 0) :
    casimirEnergy κ m₁ / casimirEnergy κ m₂ = (effRadius κ m₂ / effRadius κ m₁) \textasciicircum{} 5 := by
  rw [casimirEnergy\_\allowbreak{}ratio hκ h₁ h₂, effRadius\_\allowbreak{}div hκ h₂ h₁]
\mbox{}
/-! \#\# The geometric ansatz as an explicit structure
\mbox{}
Isolating the ansatz makes two things visible: which quantities are free, and
that the amplification ratio does not depend on any of them.
-/
\mbox{}
/-- The geometric ansatz of the framework: an internal plate geometry whose
separation is the effective radius and whose area scales as the square of the
mass. -/
structure CasimirAnsatz where
  /-- Constant of the inverse-radius law, `r\_\allowbreak{}ℓ = κ/m\_\allowbreak{}ℓ`. -/
  kappa : ℝ
  kappa\_\allowbreak{}pos : 0 \textless{} kappa
  /-- Coefficient of the plate-area ansatz, `S = areaCoeff · m²`. -/
  areaCoeff : ℝ
  areaCoeff\_\allowbreak{}pos : 0 \textless{} areaCoeff
\mbox{}
namespace CasimirAnsatz
\mbox{}
/-- Plate area of the internal geometry, `S = areaCoeff · m²`. -/
def plateArea (A : CasimirAnsatz) (m : ℝ) : ℝ := A.areaCoeff * m \textasciicircum{} 2
\mbox{}
/-- Casimir energy of the ansatz, `E = -S/d³` up to a positive prefactor;
its magnitude is `S/d³`. -/
def energy (A : CasimirAnsatz) (m : ℝ) : ℝ := A.plateArea m / (effRadius A.kappa m) \textasciicircum{} 3
\mbox{}
/-- The signed (physical, attractive) Casimir energy, `E = -C · S/d³` with `C \textgreater{} 0`. -/
def signedEnergy (A : CasimirAnsatz) (C m : ℝ) : ℝ := -C * A.energy m
\mbox{}
theorem energy\_\allowbreak{}eq\_\allowbreak{}quintic (A : CasimirAnsatz) \{m : ℝ\} (hm : m ≠ 0) :
    A.energy m = A.areaCoeff * m \textasciicircum{} 5 / A.kappa \textasciicircum{} 3 := by
  unfold energy plateArea effRadius
  field\_\allowbreak{}simp [A.kappa\_\allowbreak{}pos.ne', A.areaCoeff\_\allowbreak{}pos.ne',
    mul\_\allowbreak{}ne\_\allowbreak{}zero A.areaCoeff\_\allowbreak{}pos.ne' A.kappa\_\allowbreak{}pos.ne']
\mbox{}
/-- The amplification ratio of the ansatz is the fifth power of the mass ratio. -/
theorem energy\_\allowbreak{}ratio (A : CasimirAnsatz) \{m₁ m₂ : ℝ\} (h₁ : m₁ ≠ 0) (h₂ : m₂ ≠ 0) :
    A.energy m₁ / A.energy m₂ = (m₁ / m₂) \textasciicircum{} 5 := by
  rw [energy\_\allowbreak{}eq\_\allowbreak{}quintic A h₁, energy\_\allowbreak{}eq\_\allowbreak{}quintic A h₂]
  field\_\allowbreak{}simp [A.kappa\_\allowbreak{}pos.ne', A.areaCoeff\_\allowbreak{}pos.ne', pow\_\allowbreak{}ne\_\allowbreak{}zero 5 h₁, pow\_\allowbreak{}ne\_\allowbreak{}zero 5 h₂]
\mbox{}
/-- **Parameter-freedom of the prediction.**  Two arbitrary choices of the
framework's free constants — `κ`, the area coefficient, even the overall Casimir
prefactor — give the same amplification ratio.  The ratio `(m₁/m₂)⁵` is thus a
genuine prediction of the ansatz rather than a fit. -/
theorem energy\_\allowbreak{}ratio\_\allowbreak{}parameter\_\allowbreak{}free (A A' : CasimirAnsatz) \{m₁ m₂ : ℝ\}
    (h₁ : m₁ ≠ 0) (h₂ : m₂ ≠ 0) :
    A.energy m₁ / A.energy m₂ = A'.energy m₁ / A'.energy m₂ := by
  rw [energy\_\allowbreak{}ratio A h₁ h₂, energy\_\allowbreak{}ratio A' h₁ h₂]
\mbox{}
/-- The sign convention of the Casimir effect does not affect the amplification:
the physical energy is negative (attractive), and the ratio is unchanged. -/
theorem signedEnergy\_\allowbreak{}ratio (A : CasimirAnsatz) \{C m₁ m₂ : ℝ\} (hC : C ≠ 0)
    (h₁ : m₁ ≠ 0) (h₂ : m₂ ≠ 0) :
    A.signedEnergy C m₁ / A.signedEnergy C m₂ = (m₁ / m₂) \textasciicircum{} 5 := by
  unfold signedEnergy
  rw [neg\_\allowbreak{}mul, neg\_\allowbreak{}mul, neg\_\allowbreak{}div\_\allowbreak{}neg\_\allowbreak{}eq, mul\_\allowbreak{}div\_\allowbreak{}mul\_\allowbreak{}left \_\allowbreak{} \_\allowbreak{} hC, energy\_\allowbreak{}ratio A h₁ h₂]
\mbox{}
end CasimirAnsatz
\mbox{}
end
\mbox{}
end TheoreticalFramework
\mbox{}
}

\clearpage


\subsection*{\texttt{TheoreticalFramework/Leptons.lean}}
\vspace{2pt}
{\ttfamily\scriptsize\setlength{\parindent}{0pt}\setlength{\baselineskip}{7.6pt}\setlength{\parskip}{0pt}\raggedright\obeyspaces\obeylines
/-
\# TheoreticalFramework.Leptons
\mbox{}
Numerical instantiation of the compactness framework for the three charged
leptons, and the Casimir amplification ratios of the muonic and tauonic regimes.
\mbox{}
The lepton masses enter only through the ratios `m\_\allowbreak{}μ/m\_\allowbreak{}e ≈ 207` and
`m\_\allowbreak{}τ/m\_\allowbreak{}e ≈ 3477`.  Both integers are checked here against the tabulated masses
(`0.511`, `105.66`, `1777` MeV/c²) at the precision at which the framework
states them, and the amplification ratios are then computed exactly and bounded
in the published orders of magnitude.
-/
import TheoreticalFramework.Basic
\mbox{}
namespace TheoreticalFramework
\mbox{}
noncomputable section
\mbox{}
/-! \#\# Lepton data -/
\mbox{}
/-- Electron mass, MeV/c². -/
def massElectronMeV : ℝ := 0.511
\mbox{}
/-- Muon mass, MeV/c². -/
def massMuonMeV : ℝ := 105.66
\mbox{}
/-- Tauon mass, MeV/c². -/
def massTauonMeV : ℝ := 1777
\mbox{}
/-- Muon-to-electron mass ratio, as used throughout the framework. -/
def massRatioMuon : ℝ := 207
\mbox{}
/-- Tauon-to-electron mass ratio, as used throughout the framework. -/
def massRatioTauon : ℝ := 3477
\mbox{}
theorem massRatioMuon\_\allowbreak{}pos : 0 \textless{} massRatioMuon := by
  norm\_\allowbreak{}num [massRatioMuon]
\mbox{}
theorem massRatioTauon\_\allowbreak{}pos : 0 \textless{} massRatioTauon := by
  norm\_\allowbreak{}num [massRatioTauon]
\mbox{}
/-- The integer ratio `207` reproduces the tabulated muon/electron mass ratio to
better than a quarter of a unit. -/
theorem massRatioMuon\_\allowbreak{}consistent :
    |massMuonMeV / massElectronMeV - massRatioMuon| ≤ 0.25 := by
  unfold massMuonMeV massElectronMeV massRatioMuon
  rw [abs\_\allowbreak{}le]
  constructor \textless{};\textgreater{} norm\_\allowbreak{}num
\mbox{}
/-- The integer ratio `3477` reproduces the tabulated tauon/electron mass ratio
to better than half a unit. -/
theorem massRatioTauon\_\allowbreak{}consistent :
    |massTauonMeV / massElectronMeV - massRatioTauon| ≤ 0.5 := by
  unfold massTauonMeV massElectronMeV massRatioTauon
  rw [abs\_\allowbreak{}le]
  constructor \textless{};\textgreater{} norm\_\allowbreak{}num
\mbox{}
/-! \#\# Casimir amplification, muon and tauon versus electron
\mbox{}
The general ratio `casimirEnergy\_\allowbreak{}ratio` is instantiated with masses
`m\_\allowbreak{}ℓ = (m\_\allowbreak{}ℓ/m\_\allowbreak{}e) · m\_\allowbreak{}e`, so the result is a statement about the physics, not
about a bare numeral.  The numeral is then evaluated exactly.
-/
\mbox{}
/-- **Muonic Casimir amplification**, in the form `E\_\allowbreak{}μ / E\_\allowbreak{}e = (m\_\allowbreak{}μ/m\_\allowbreak{}e)⁵`. -/
theorem casimir\_\allowbreak{}amplification\_\allowbreak{}muon\_\allowbreak{}general \{κ m₀ : ℝ\} (hκ : κ ≠ 0) (hm : 0 \textless{} m₀) :
    casimirEnergy κ (massRatioMuon * m₀) / casimirEnergy κ m₀ = massRatioMuon \textasciicircum{} 5 := by
  rw [casimirEnergy\_\allowbreak{}ratio hκ (mul\_\allowbreak{}ne\_\allowbreak{}zero massRatioMuon\_\allowbreak{}pos.ne' hm.ne') hm.ne']
  field\_\allowbreak{}simp
\mbox{}
/-- **Tauonic Casimir amplification**, in the form `E\_\allowbreak{}τ / E\_\allowbreak{}e = (m\_\allowbreak{}τ/m\_\allowbreak{}e)⁵`. -/
theorem casimir\_\allowbreak{}amplification\_\allowbreak{}tauon\_\allowbreak{}general \{κ m₀ : ℝ\} (hκ : κ ≠ 0) (hm : 0 \textless{} m₀) :
    casimirEnergy κ (massRatioTauon * m₀) / casimirEnergy κ m₀ = massRatioTauon \textasciicircum{} 5 := by
  rw [casimirEnergy\_\allowbreak{}ratio hκ (mul\_\allowbreak{}ne\_\allowbreak{}zero massRatioTauon\_\allowbreak{}pos.ne' hm.ne') hm.ne']
  field\_\allowbreak{}simp
\mbox{}
/-- Exact value of the muonic amplification, `207⁵ = 380 059 617 807`. -/
theorem massRatioMuon\_\allowbreak{}pow\_\allowbreak{}five : massRatioMuon \textasciicircum{} 5 = 380059617807 := by
  norm\_\allowbreak{}num [massRatioMuon]
\mbox{}
/-- Exact value of the tauonic amplification, `3477⁵ = 508 186 885 683 281 157`. -/
theorem massRatioTauon\_\allowbreak{}pow\_\allowbreak{}five : massRatioTauon \textasciicircum{} 5 = 508186885683281157 := by
  norm\_\allowbreak{}num [massRatioTauon]
\mbox{}
/-- `(207)⁵ ≈ 4 × 10¹¹`: the amplification lies between `3.8 × 10¹¹` and `4 × 10¹¹`. -/
theorem casimir\_\allowbreak{}amplification\_\allowbreak{}muon\_\allowbreak{}order :
    (3.8 : ℝ) * 10 \textasciicircum{} 11 ≤ massRatioMuon \textasciicircum{} 5 ∧ massRatioMuon \textasciicircum{} 5 ≤ 4 * 10 \textasciicircum{} 11 := by
  rw [massRatioMuon\_\allowbreak{}pow\_\allowbreak{}five]
  constructor \textless{};\textgreater{} norm\_\allowbreak{}num
\mbox{}
/-- `(3477)⁵ ≈ 5.1 × 10¹⁷`: the amplification lies between `5.08 × 10¹⁷` and `5.1 × 10¹⁷`. -/
theorem casimir\_\allowbreak{}amplification\_\allowbreak{}tauon\_\allowbreak{}order :
    (5.08 : ℝ) * 10 \textasciicircum{} 17 ≤ massRatioTauon \textasciicircum{} 5 ∧ massRatioTauon \textasciicircum{} 5 ≤ (5.1 : ℝ) * 10 \textasciicircum{} 17 := by
  rw [massRatioTauon\_\allowbreak{}pow\_\allowbreak{}five]
  constructor \textless{};\textgreater{} norm\_\allowbreak{}num
\mbox{}
/-- The tauon amplifies the vacuum coupling by a further factor
`(3477/207)⁵ ≈ 1.3 × 10⁶` over the muon. -/
theorem casimir\_\allowbreak{}amplification\_\allowbreak{}tauon\_\allowbreak{}over\_\allowbreak{}muon :
    (1.3 : ℝ) * 10 \textasciicircum{} 6 ≤ (massRatioTauon / massRatioMuon) \textasciicircum{} 5 ∧
      (massRatioTauon / massRatioMuon) \textasciicircum{} 5 ≤ (1.34 : ℝ) * 10 \textasciicircum{} 6 := by
  unfold massRatioTauon massRatioMuon
  constructor \textless{};\textgreater{} norm\_\allowbreak{}num
\mbox{}
/-- The amplification is monotone in the mass ratio: the lepton hierarchy
`e \textless{} μ \textless{} τ` is a strict hierarchy of vacuum couplings. -/
theorem casimir\_\allowbreak{}amplification\_\allowbreak{}hierarchy :
    (1 : ℝ) \textless{} massRatioMuon \textasciicircum{} 5 ∧ massRatioMuon \textasciicircum{} 5 \textless{} massRatioTauon \textasciicircum{} 5 := by
  rw [massRatioMuon\_\allowbreak{}pow\_\allowbreak{}five, massRatioTauon\_\allowbreak{}pow\_\allowbreak{}five]
  constructor \textless{};\textgreater{} norm\_\allowbreak{}num
\mbox{}
/-! \#\# Equivalent magnetic field, and where the exponent `3` comes from
\mbox{}
The source text associates the Casimir energy with an equivalent magnetic field
through `B² $\propto$ E\_\allowbreak{}C / Volume`.  Two volume ansätze are possible, and they give
different exponents.  Both are proved here so that the framework's use of
`B/B₀ $\propto$ m³` is tied to a stated assumption rather than left implicit.
\mbox{}
* With the source's volume ansatz `V $\propto$ S·d³` one gets `B² $\propto$ m⁶`, i.e. `B $\propto$ m³`.
* With the geometrically literal plate volume `V = S·d` one gets `B² $\propto$ m⁴`, i.e. `B $\propto$ m²`.
\mbox{}
The first is the exponent the framework uses, and it coincides with the velocity
exponent `n - 2 = 3`; the second is recorded because it is what a literal plate
volume gives, and the discrepancy is a property of the ansatz, not of the
algebra.
-/
\mbox{}
/-- Equivalent-field amplification under the source's volume ansatz `V $\propto$ S·d³`:
`B $\propto$ m³`. -/
def fieldAmplificationSource (m₁ m₂ : ℝ) : ℝ := (m₁ / m₂) \textasciicircum{} 3
\mbox{}
/-- Equivalent-field amplification under the literal plate volume `V = S·d`:
`B $\propto$ m²`. -/
def fieldAmplificationPlate (m₁ m₂ : ℝ) : ℝ := (m₁ / m₂) \textasciicircum{} 2
\mbox{}
/-- Source ansatz: `B² $\propto$ E\_\allowbreak{}C / (S·d³) $\propto$ m⁶`, hence `B $\propto$ m³`. -/
theorem field\_\allowbreak{}amplification\_\allowbreak{}source (κ m₁ m₂ : ℝ) (hκ : 0 \textless{} κ) (h₁ : 0 \textless{} m₁) (h₂ : 0 \textless{} m₂) :
    (casimirEnergy κ m₁ / (m₁ \textasciicircum{} 2 * effRadius κ m₁ \textasciicircum{} 3)) /
        (casimirEnergy κ m₂ / (m₂ \textasciicircum{} 2 * effRadius κ m₂ \textasciicircum{} 3)) = (m₁ / m₂) \textasciicircum{} 6 := by
  rw [casimirEnergy\_\allowbreak{}eq\_\allowbreak{}quintic hκ.ne' h₁.ne', casimirEnergy\_\allowbreak{}eq\_\allowbreak{}quintic hκ.ne' h₂.ne']
  unfold effRadius
  field\_\allowbreak{}simp
\mbox{}
/-- Literal plate volume: `B² $\propto$ E\_\allowbreak{}C / (S·d) $\propto$ m⁴`, hence `B $\propto$ m²`. -/
theorem field\_\allowbreak{}amplification\_\allowbreak{}plate (κ m₁ m₂ : ℝ) (hκ : 0 \textless{} κ) (h₁ : 0 \textless{} m₁) (h₂ : 0 \textless{} m₂) :
    (casimirEnergy κ m₁ / (m₁ \textasciicircum{} 2 * effRadius κ m₁)) /
        (casimirEnergy κ m₂ / (m₂ \textasciicircum{} 2 * effRadius κ m₂)) = (m₁ / m₂) \textasciicircum{} 4 := by
  rw [casimirEnergy\_\allowbreak{}eq\_\allowbreak{}quintic hκ.ne' h₁.ne', casimirEnergy\_\allowbreak{}eq\_\allowbreak{}quintic hκ.ne' h₂.ne']
  unfold effRadius
  field\_\allowbreak{}simp
\mbox{}
/-- The source's field amplification for the muon, `(207)³ ≈ 9 × 10⁶`. -/
theorem field\_\allowbreak{}amplification\_\allowbreak{}muon\_\allowbreak{}order :
    (8.8 : ℝ) * 10 \textasciicircum{} 6 ≤ fieldAmplificationSource massRatioMuon 1 ∧
      fieldAmplificationSource massRatioMuon 1 ≤ (8.9 : ℝ) * 10 \textasciicircum{} 6 := by
  unfold fieldAmplificationSource massRatioMuon
  constructor \textless{};\textgreater{} norm\_\allowbreak{}num
\mbox{}
end
\mbox{}
end TheoreticalFramework
\mbox{}
}

\clearpage


\subsection*{\texttt{TheoreticalFramework/Dimensions.lean}}
\vspace{2pt}
{\ttfamily\scriptsize\setlength{\parindent}{0pt}\setlength{\baselineskip}{7.6pt}\setlength{\parskip}{0pt}\raggedright\obeyspaces\obeylines
/-
\# TheoreticalFramework.Dimensions
\mbox{}
The generalized `n`-dimensional kinematic framework.
\mbox{}
The framework's generalized dimension is not a free parameter: it is read off the
quintic mass dependence of the anomalous-moment contribution,
\mbox{}
  `a\_\allowbreak{}ℓ = a\_\allowbreak{}QED + a\_\allowbreak{}EW+had(m\_\allowbreak{}ℓ⁵) + a\_\allowbreak{}str(m\_\allowbreak{}ℓ, G\_\allowbreak{}ℓ)`,
\mbox{}
whose structural term scales as `m⁵` (`casimirEnergy\_\allowbreak{}eq\_\allowbreak{}quintic`).  Setting
`n = 5` makes the velocity exponent `n − 2 = 3`, which is the exponent that
appears in the ejection-velocity law and, independently, in the equivalent
magnetic-field amplification of `TheoreticalFramework.Leptons`.
\mbox{}
The velocity law itself is stated for arbitrary `n ≥ 2` and then specialised:
\mbox{}
  `vf² = v₂² · (m\_\allowbreak{}ℓ/m₀)\textasciicircum{}(n−2)`.
\mbox{}
`vf` is defined as the non-negative root of that equation, and the paper's
multiplicative form `vf = v₂ · √((m\_\allowbreak{}ℓ/m₀)\textasciicircum{}(n−2))` is proved equivalent for
`v₂ ≥ 0`.
\mbox{}
The generalized units of the framework, `[Cₙ] = [vₙ] = m\textasciicircum{}(2/n)/s`, enter only
through exponent bookkeeping; that bookkeeping is proved in
`generalised\_\allowbreak{}unit\_\allowbreak{}exponents` and `time\_\allowbreak{}exponent\_\allowbreak{}balance`.
-/
import TheoreticalFramework.Leptons
\mbox{}
namespace TheoreticalFramework
\mbox{}
noncomputable section
\mbox{}
/-! \#\# The generalized dimension -/
\mbox{}
/-- Exponent of the mass in the Casimir energy of the compactness ansatz. -/
def casimirMassExponent : ℕ := 5
\mbox{}
/-- The generalized dimension, fixed by the quintic mass dependence: `n = 5`. -/
def generalizedDimension : ℕ := casimirMassExponent
\mbox{}
/-- The generalized dimension is five. -/
theorem generalizedDimension\_\allowbreak{}eq\_\allowbreak{}five : generalizedDimension = 5 := rfl
\mbox{}
/-! \#\# The anomalous moment whose quintic term fixes `n`
\mbox{}
The framework decomposes the anomalous moment of a lepton into a QED part, an
electroweak-plus-hadronic part and a structural part, the last being quintic in
the mass.  The quintic exponent is what fixes `n = 5`; it is not an independent
input. -/
\mbox{}
/-- The structural part of the anomalous moment, `a\_\allowbreak{}str = C · m\textasciicircum{}n`. -/
def structuralTerm (C m : ℝ) : ℝ := C * m \textasciicircum{} generalizedDimension
\mbox{}
/-- The framework's decomposition `a\_\allowbreak{}ℓ = a\_\allowbreak{}QED + a\_\allowbreak{}EW+had + a\_\allowbreak{}str`. -/
def anomalousMoment (aQED aEWHad C m : ℝ) : ℝ := aQED + aEWHad + structuralTerm C m
\mbox{}
/-- The structural term is quintic in the mass. -/
theorem structuralTerm\_\allowbreak{}quintic (C m : ℝ) : structuralTerm C m = C * (m : ℝ) \textasciicircum{} 5 := by
  unfold structuralTerm
  rw [generalizedDimension\_\allowbreak{}eq\_\allowbreak{}five]
\mbox{}
/-- The structural term amplifies between two leptons by the fifth power of the
mass ratio: the same number as the Casimir amplification, which is the reason a
single exponent serves both. -/
theorem structuralTerm\_\allowbreak{}ratio (C m₁ m₂ : ℝ) (hC : C ≠ 0) (h₁ : m₁ ≠ 0) (h₂ : m₂ ≠ 0) :
    structuralTerm C m₁ / structuralTerm C m₂ = (m₁ / m₂) \textasciicircum{} 5 := by
  unfold structuralTerm
  rw [generalizedDimension\_\allowbreak{}eq\_\allowbreak{}five, div\_\allowbreak{}pow]
  field\_\allowbreak{}simp [hC, pow\_\allowbreak{}ne\_\allowbreak{}zero 5 h₁, pow\_\allowbreak{}ne\_\allowbreak{}zero 5 h₂]
\mbox{}
/-- Exponent of the mass ratio in the velocity law, `n − 2`. -/
def velocityExponent : ℕ := generalizedDimension - 2
\mbox{}
/-- With `n = 5` the velocity law carries the cube of the mass ratio. -/
theorem velocityExponent\_\allowbreak{}eq\_\allowbreak{}three : velocityExponent = 3 := rfl
\mbox{}
/-- The velocity exponent is `n − 2` for the framework's `n`. -/
theorem velocityExponent\_\allowbreak{}eq (n : ℕ) (hn : n = generalizedDimension) :
    velocityExponent = n - 2 := by
  subst hn
  rfl
\mbox{}
/-! \#\# Generalized units
\mbox{}
The framework assigns to the generalized velocity the units `m\textasciicircum{}(2/n)/s`, so that
the `n`-th power carries `m²/sⁿ`; multiplying by `dt₀\textasciicircum{}(n−2)` then carries
`m²/s²`, the units of `v₂²`.  Only the exponents matter, and they are exact.
-/
\mbox{}
/-- Unit bookkeeping of `[vₙ] = m\textasciicircum{}(2/n)/s`: the mass exponent balances to `2`. -/
theorem generalised\_\allowbreak{}unit\_\allowbreak{}exponents (n : ℚ) (hn : n ≠ 0) : n * (2 / n) = 2 := by
  field\_\allowbreak{}simp
\mbox{}
/-- Unit bookkeeping of the time factor `dt₀\textasciicircum{}(n−2)`: the exponents close on `n`. -/
theorem time\_\allowbreak{}exponent\_\allowbreak{}balance (n : ℚ) : (n - 2) + 2 = n := by
  ring
\mbox{}
/-- The generalized-velocity relation `vₙⁿ · dt₀\textasciicircum{}(n−2) = v₂²` of the framework. -/
def GeneralisedVelocityRelation (n : ℕ) (vn dt0 v2 : ℝ) : Prop :=
  vn \textasciicircum{} n * dt0 \textasciicircum{} (n - 2) = v2 \textasciicircum{} 2
\mbox{}
/-- **Scale substitution.**  Replacing the time scale by `q·dt₀` multiplies the
relation by `q\textasciicircum{}(n−2)`, which is the framework's substitution rule
`vₙ·(q·dt₀)\textasciicircum{}(n−2) = v₂²·q\textasciicircum{}(n−2)`; the parameter `q = m\_\allowbreak{}ℓ/m₀` is what links the
muonic and tauonic regimes. -/
theorem scale\_\allowbreak{}substitution (n : ℕ) \{vn dt0 q v2 : ℝ\}
    (h : vn * dt0 \textasciicircum{} (n - 2) = v2 \textasciicircum{} 2) :
    vn * (q * dt0) \textasciicircum{} (n - 2) = v2 \textasciicircum{} 2 * q \textasciicircum{} (n - 2) := by
  rw [mul\_\allowbreak{}pow, ← h]
  ring
\mbox{}
/-! \#\# The velocity law -/
\mbox{}
/-- Ejection velocity of particles produced in lepton–lepton collisions in the
generalized framework: the non-negative root of `vf² = v₂²·r\textasciicircum{}(n−2)`, where `r`
is the mass ratio `m\_\allowbreak{}ℓ/m₀`. -/
def generalisedVelocity (n : ℕ) (v₂ r : ℝ) : ℝ := Real.sqrt (v₂ \textasciicircum{} 2 * r \textasciicircum{} (n - 2))
\mbox{}
/-- **The velocity law**, for arbitrary generalized dimension `n`:
`vf² = v₂² · (m\_\allowbreak{}ℓ/m₀)\textasciicircum{}(n−2)`. -/
theorem velocity\_\allowbreak{}law (n : ℕ) (v₂ r : ℝ) (hr : 0 ≤ r) :
    generalisedVelocity n v₂ r \textasciicircum{} 2 = v₂ \textasciicircum{} 2 * r \textasciicircum{} (n - 2) := by
  unfold generalisedVelocity
  exact Real.sq\_\allowbreak{}sqrt (mul\_\allowbreak{}nonneg (sq\_\allowbreak{}nonneg v₂) (pow\_\allowbreak{}nonneg hr \_\allowbreak{}))
\mbox{}
/-- The paper's form of the law, `vf = v₂ · √((m\_\allowbreak{}ℓ/m₀)\textasciicircum{}(n−2))`, agrees with the
definition for non-negative input velocity. -/
theorem generalisedVelocity\_\allowbreak{}eq\_\allowbreak{}mul (n : ℕ) \{v₂ r : ℝ\} (hv : 0 ≤ v₂) :
    generalisedVelocity n v₂ r = v₂ * Real.sqrt (r \textasciicircum{} (n - 2)) := by
  unfold generalisedVelocity
  rw [Real.sqrt\_\allowbreak{}mul (sq\_\allowbreak{}nonneg v₂), Real.sqrt\_\allowbreak{}sq hv]
\mbox{}
/-- Ejection velocities are non-negative. -/
theorem generalisedVelocity\_\allowbreak{}nonneg (n : ℕ) (v₂ r : ℝ) : 0 ≤ generalisedVelocity n v₂ r :=
  Real.sqrt\_\allowbreak{}nonneg \_\allowbreak{}
\mbox{}
/-- Two-sided bounds on an ejection velocity, from exact bounds on `vf²`.
This is the workhorse that turns a computed square into a velocity interval
without any rounding step. -/
theorem generalisedVelocity\_\allowbreak{}bounds \{n : ℕ\} \{v₂ r a b : ℝ\} (hb : 0 ≤ b)
    (h₁ : a \textasciicircum{} 2 ≤ v₂ \textasciicircum{} 2 * r \textasciicircum{} (n - 2)) (h₂ : v₂ \textasciicircum{} 2 * r \textasciicircum{} (n - 2) ≤ b \textasciicircum{} 2) :
    a ≤ generalisedVelocity n v₂ r ∧ generalisedVelocity n v₂ r ≤ b := by
  unfold generalisedVelocity
  exact ⟨Real.le\_\allowbreak{}sqrt\_\allowbreak{}of\_\allowbreak{}sq\_\allowbreak{}le h₁, (Real.sqrt\_\allowbreak{}le\_\allowbreak{}left hb).2 h₂⟩
\mbox{}
/-- The velocity law at the framework's generalized dimension:
`vf² = v₂² · (m\_\allowbreak{}ℓ/m₀)³`. -/
theorem velocity\_\allowbreak{}law\_\allowbreak{}five (v₂ r : ℝ) (hr : 0 ≤ r) :
    generalisedVelocity generalizedDimension v₂ r \textasciicircum{} 2 = v₂ \textasciicircum{} 2 * r \textasciicircum{} 3 := by
  rw [velocity\_\allowbreak{}law generalizedDimension v₂ r hr, generalizedDimension\_\allowbreak{}eq\_\allowbreak{}five]
\mbox{}
/-- The law's explicit linear-in-`v₂` form at `n = 5`: `vf = v₂·√(r³)`. -/
theorem generalisedVelocity\_\allowbreak{}five\_\allowbreak{}eq\_\allowbreak{}mul \{v₂ r : ℝ\} (hv : 0 ≤ v₂) :
    generalisedVelocity generalizedDimension v₂ r = v₂ * Real.sqrt (r \textasciicircum{} 3) := by
  rw [generalisedVelocity\_\allowbreak{}eq\_\allowbreak{}mul generalizedDimension hv, generalizedDimension\_\allowbreak{}eq\_\allowbreak{}five]
\mbox{}
/-- The ejection velocity grows with the lepton mass ratio: heavier collision
partners eject faster. -/
theorem generalisedVelocity\_\allowbreak{}mono \{n : ℕ\} \{v₂ r₁ r₂ : ℝ\} (hr₁ : 0 ≤ r₁) (h : r₁ ≤ r₂) :
    generalisedVelocity n v₂ r₁ ≤ generalisedVelocity n v₂ r₂ := by
  unfold generalisedVelocity
  apply Real.sqrt\_\allowbreak{}le\_\allowbreak{}sqrt
  exact mul\_\allowbreak{}le\_\allowbreak{}mul\_\allowbreak{}of\_\allowbreak{}nonneg\_\allowbreak{}left (pow\_\allowbreak{}le\_\allowbreak{}pow\_\allowbreak{}left₀ hr₁ h \_\allowbreak{}) (sq\_\allowbreak{}nonneg v₂)
\mbox{}
/-- The ejection velocity grows with the collision velocity. -/
theorem generalisedVelocity\_\allowbreak{}mono\_\allowbreak{}v₂ \{n : ℕ\} \{v₂ v₂' r : ℝ\} (hv : 0 ≤ v₂) (h : v₂ ≤ v₂')
    (hr : 0 ≤ r) :
    generalisedVelocity n v₂ r ≤ generalisedVelocity n v₂' r := by
  unfold generalisedVelocity
  apply Real.sqrt\_\allowbreak{}le\_\allowbreak{}sqrt
  exact mul\_\allowbreak{}le\_\allowbreak{}mul\_\allowbreak{}of\_\allowbreak{}nonneg\_\allowbreak{}right (pow\_\allowbreak{}le\_\allowbreak{}pow\_\allowbreak{}left₀ hv h 2) (pow\_\allowbreak{}nonneg hr \_\allowbreak{})
\mbox{}
/-- The ejection velocity vanishes with the collision velocity: no collisions,
no propulsion. -/
theorem generalisedVelocity\_\allowbreak{}zero (n : ℕ) (r : ℝ) : generalisedVelocity n 0 r = 0 := by
  unfold generalisedVelocity
  simp
\mbox{}
end
\mbox{}
end TheoreticalFramework
\mbox{}
}

\clearpage


\subsection*{\texttt{TheoreticalFramework/Velocity.lean}}
\vspace{2pt}
{\ttfamily\scriptsize\setlength{\parindent}{0pt}\setlength{\baselineskip}{7.6pt}\setlength{\parskip}{0pt}\raggedright\obeyspaces\obeylines
/-
\# TheoreticalFramework.Velocity
\mbox{}
Velocity scaling law, amplification factors and ejection velocities of the
muonic and tauonic regimes, with every physical constant bounded rigorously.
\mbox{}
Units.  The framework states collision velocities in km/h and ejection
velocities in m/s and km/s.  All conversions are proved; no numeric value is
asserted without a surrounding interval, and each interval is closed by an exact
computation of `vf²` together with `generalisedVelocity\_\allowbreak{}bounds`.
\mbox{}
The two amplification factors of the framework are
`(m\_\allowbreak{}μ/m\_\allowbreak{}e)\textasciicircum{}(n−2) = 207³ = 8 869 743 ≈ 8.9 × 10⁶` and
`(m\_\allowbreak{}τ/m\_\allowbreak{}e)\textasciicircum{}(n−2) = 3477³ = 42 035 292 333 ≈ 4.2 × 10¹⁰`.
\mbox{}
The muonic ceiling reproduces the source's `≈ 20 700 km/s`.  The tauonic ceiling
does *not* reproduce the source's `1.34 × 10¹² m/s`: the exact law gives
`4.2403 × 10¹²`, and the discrepancy is localised in
`source\_\allowbreak{}tauon\_\allowbreak{}max\_\allowbreak{}velocity\_\allowbreak{}underestimates` — the source evaluates
`√(1.80 × 10²⁵)` as `1.34 × 10¹²` where it is `4.24 × 10¹²`, a `√10` slip in the
decimal exponent.  The published figure is kept as a separate definition so that
the two can be compared without either being hidden.
-/
import TheoreticalFramework.Dimensions
\mbox{}
namespace TheoreticalFramework
\mbox{}
noncomputable section
\mbox{}
/-! \#\# Units and framework constants -/
\mbox{}
/-- Kilometres per hour to metres per second, `1 km/h = 5/18 m/s`. -/
def kmhToMs (v : ℝ) : ℝ := v * 5 / 18
\mbox{}
/-- The conversion agrees with the definitional `1000/3600`. -/
theorem kmhToMs\_\allowbreak{}eq (v : ℝ) : kmhToMs v = v * 1000 / 3600 := by
  unfold kmhToMs
  ring
\mbox{}
/-- Standard gravity, `g₀ = 9.80665 m/s²`, used to express specific impulse. -/
def standardGravity : ℝ := 9.80665
\mbox{}
theorem standardGravity\_\allowbreak{}pos : 0 \textless{} standardGravity := by
  norm\_\allowbreak{}num [standardGravity]
\mbox{}
/-- Speed of light in vacuum, m/s. -/
def speedOfLight : ℝ := 299792458
\mbox{}
/-- Specific impulse, `Isp = v\_\allowbreak{}e / g₀`, in seconds. -/
def specificImpulse (ve : ℝ) : ℝ := ve / standardGravity
\mbox{}
/-- Vacuum specific impulse of a Raptor engine, seconds: the chemical benchmark
of the source text. -/
def raptorSpecificImpulse : ℝ := 363
\mbox{}
/-! \#\# Amplification factors -/
\mbox{}
/-- Upper end of the muonic input window, `v₂ = 25 000 km/h`. -/
def collisionVelocityMuonMax : ℝ := kmhToMs 25000
\mbox{}
/-- The muonic input ceiling in m/s. -/
theorem collisionVelocityMuonMax\_\allowbreak{}eq : collisionVelocityMuonMax = 62500 / 9 := by
  unfold collisionVelocityMuonMax kmhToMs
  norm\_\allowbreak{}num
\mbox{}
/-- `25 000 km/h ≈ 6 944 m/s`. -/
theorem collisionVelocityMuonMax\_\allowbreak{}bounds :
    6944 ≤ collisionVelocityMuonMax ∧ collisionVelocityMuonMax ≤ 6945 := by
  rw [collisionVelocityMuonMax\_\allowbreak{}eq]
  constructor \textless{};\textgreater{} norm\_\allowbreak{}num
\mbox{}
/-- Muonic amplification factor, `(m\_\allowbreak{}μ/m\_\allowbreak{}e)\textasciicircum{}(n−2)`. -/
def muonAmplification : ℝ := massRatioMuon \textasciicircum{} velocityExponent
\mbox{}
/-- Tauonic amplification factor, `(m\_\allowbreak{}τ/m\_\allowbreak{}e)\textasciicircum{}(n−2)`. -/
def tauonAmplification : ℝ := massRatioTauon \textasciicircum{} velocityExponent
\mbox{}
/-- `207³ = 8 869 743`. -/
theorem muonAmplification\_\allowbreak{}eq : muonAmplification = 8869743 := by
  unfold muonAmplification velocityExponent generalizedDimension casimirMassExponent
  norm\_\allowbreak{}num [massRatioMuon]
\mbox{}
/-- `3477³ = 42 035 292 333`. -/
theorem tauonAmplification\_\allowbreak{}eq : tauonAmplification = 42035292333 := by
  unfold tauonAmplification velocityExponent generalizedDimension casimirMassExponent
  norm\_\allowbreak{}num [massRatioTauon]
\mbox{}
/-- The muonic amplification is `(207)³ ≈ 8.9 × 10⁶`. -/
theorem muonAmplification\_\allowbreak{}order :
    (8.8 : ℝ) * 10 \textasciicircum{} 6 ≤ muonAmplification ∧ muonAmplification ≤ (8.9 : ℝ) * 10 \textasciicircum{} 6 := by
  rw [muonAmplification\_\allowbreak{}eq]
  constructor \textless{};\textgreater{} norm\_\allowbreak{}num
\mbox{}
/-- The tauonic amplification is `(3477)³ ≈ 4.2 × 10¹⁰`. -/
theorem tauonAmplification\_\allowbreak{}order :
    (4.2 : ℝ) * 10 \textasciicircum{} 10 ≤ tauonAmplification ∧ tauonAmplification ≤ (4.21 : ℝ) * 10 \textasciicircum{} 10 := by
  rw [tauonAmplification\_\allowbreak{}eq]
  constructor \textless{};\textgreater{} norm\_\allowbreak{}num
\mbox{}
/-- The tauonic amplification exceeds the muonic one by `(3477/207)³ ≈ 4.74 × 10³`. -/
theorem tauon\_\allowbreak{}over\_\allowbreak{}muon\_\allowbreak{}amplification :
    (4.73 : ℝ) * 10 \textasciicircum{} 3 ≤ tauonAmplification / muonAmplification ∧
      tauonAmplification / muonAmplification ≤ (4.74 : ℝ) * 10 \textasciicircum{} 3 := by
  rw [tauonAmplification\_\allowbreak{}eq, muonAmplification\_\allowbreak{}eq]
  constructor
  · rw [le\_\allowbreak{}div\_\allowbreak{}iff₀ (by norm\_\allowbreak{}num)]
    norm\_\allowbreak{}num
  · rw [div\_\allowbreak{}le\_\allowbreak{}iff₀ (by norm\_\allowbreak{}num)]
    norm\_\allowbreak{}num
\mbox{}
/-! \#\# Muonic regime: ejection velocity -/
\mbox{}
/-- Maximum ejection velocity of the muonic regime, at `v₂ = 25 000 km/h`. -/
def muonMaxEjection : ℝ :=
  generalisedVelocity generalizedDimension collisionVelocityMuonMax massRatioMuon
\mbox{}
/-- The ejection velocity obeys the law: `vf² = v₂² · 207³`. -/
theorem muonMaxEjection\_\allowbreak{}sq :
    muonMaxEjection \textasciicircum{} 2 = collisionVelocityMuonMax \textasciicircum{} 2 * massRatioMuon \textasciicircum{} 3 := by
  unfold muonMaxEjection
  exact velocity\_\allowbreak{}law\_\allowbreak{}five \_\allowbreak{} \_\allowbreak{} massRatioMuon\_\allowbreak{}pos.le
\mbox{}
/-- `vf ≈ 2.06820 × 10⁷ m/s` for the muonic ceiling. -/
theorem muonMaxEjection\_\allowbreak{}bounds :
    (2.06820 : ℝ) * 10 \textasciicircum{} 7 ≤ muonMaxEjection ∧
      muonMaxEjection ≤ (2.06821 : ℝ) * 10 \textasciicircum{} 7 := by
  have h₁ : ((2.06820 : ℝ) * 10 \textasciicircum{} 7) \textasciicircum{} 2 ≤
      collisionVelocityMuonMax \textasciicircum{} 2 * massRatioMuon \textasciicircum{} 3 := by
    rw [collisionVelocityMuonMax\_\allowbreak{}eq]
    norm\_\allowbreak{}num [massRatioMuon]
  have h₂ : collisionVelocityMuonMax \textasciicircum{} 2 * massRatioMuon \textasciicircum{} 3 ≤
      ((2.06821 : ℝ) * 10 \textasciicircum{} 7) \textasciicircum{} 2 := by
    rw [collisionVelocityMuonMax\_\allowbreak{}eq]
    norm\_\allowbreak{}num [massRatioMuon]
  have := generalisedVelocity\_\allowbreak{}bounds (n := generalizedDimension)
    (v₂ := collisionVelocityMuonMax) (r := massRatioMuon)
    (a := (2.06820 : ℝ) * 10 \textasciicircum{} 7) (b := (2.06821 : ℝ) * 10 \textasciicircum{} 7)
    (show (0 : ℝ) ≤ (2.06821 : ℝ) * 10 \textasciicircum{} 7 by norm\_\allowbreak{}num) h₁ h₂
  simpa [muonMaxEjection] using this
\mbox{}
/-- The source's `≈ 20 700 km/s`, proved as an interval in km/s. -/
theorem muonMaxEjection\_\allowbreak{}km\_\allowbreak{}per\_\allowbreak{}s :
    (2.068 : ℝ) * 10 \textasciicircum{} 4 ≤ muonMaxEjection / 1000 ∧
      muonMaxEjection / 1000 ≤ (2.07 : ℝ) * 10 \textasciicircum{} 4 := by
  obtain ⟨h₁, h₂⟩ := muonMaxEjection\_\allowbreak{}bounds
  constructor
  · rw [le\_\allowbreak{}div\_\allowbreak{}iff₀ (by norm\_\allowbreak{}num)]
    nlinarith
  · rw [div\_\allowbreak{}le\_\allowbreak{}iff₀ (by norm\_\allowbreak{}num)]
    nlinarith
\mbox{}
/-- The muonic ejection velocity is subluminal: the muonic regime is compatible
with ordinary special relativity. -/
theorem muonMaxEjection\_\allowbreak{}subluminal : muonMaxEjection \textless{} speedOfLight := by
  have := muonMaxEjection\_\allowbreak{}bounds.2
  unfold speedOfLight
  linarith
\mbox{}
/-- The muonic input window is subluminal as well. -/
theorem collisionVelocityMuonMax\_\allowbreak{}subluminal : collisionVelocityMuonMax \textless{} speedOfLight := by
  have := collisionVelocityMuonMax\_\allowbreak{}bounds.2
  unfold speedOfLight
  linarith
\mbox{}
/-- Specific impulse of the muonic ceiling, `Isp = v\_\allowbreak{}e/g₀ ≈ 2.109 × 10⁶ s`. -/
theorem muonSpecificImpulse\_\allowbreak{}bounds :
    (2.108 : ℝ) * 10 \textasciicircum{} 6 ≤ specificImpulse muonMaxEjection ∧
      specificImpulse muonMaxEjection ≤ (2.109 : ℝ) * 10 \textasciicircum{} 6 := by
  obtain ⟨h₁, h₂⟩ := muonMaxEjection\_\allowbreak{}bounds
  have hg := standardGravity\_\allowbreak{}pos
  constructor
  · rw [specificImpulse, le\_\allowbreak{}div\_\allowbreak{}iff₀ hg]
    unfold standardGravity at hg ⊢
    nlinarith
  · rw [specificImpulse, div\_\allowbreak{}le\_\allowbreak{}iff₀ hg]
    unfold standardGravity at hg ⊢
    nlinarith
\mbox{}
/-- The muonic impulse exceeds the Raptor benchmark by a factor of about
`5.8 × 10³`, i.e. "roughly 6 000 times" as the source states. -/
theorem muonImpulse\_\allowbreak{}over\_\allowbreak{}raptor :
    (5.80 : ℝ) * 10 \textasciicircum{} 3 ≤ specificImpulse muonMaxEjection / raptorSpecificImpulse ∧
      specificImpulse muonMaxEjection / raptorSpecificImpulse ≤ (5.82 : ℝ) * 10 \textasciicircum{} 3 := by
  obtain ⟨h₁, h₂⟩ := muonSpecificImpulse\_\allowbreak{}bounds
  unfold raptorSpecificImpulse
  constructor
  · rw [le\_\allowbreak{}div\_\allowbreak{}iff₀ (by norm\_\allowbreak{}num)]
    nlinarith
  · rw [div\_\allowbreak{}le\_\allowbreak{}iff₀ (by norm\_\allowbreak{}num)]
    nlinarith
\mbox{}
/-! \#\# Tauonic regime: ejection velocity -/
\mbox{}
/-- Input ceiling of the tauonic regime: the muonic regime's maximum ejection
velocity, `≈ 20 682 km/s`.  Taking the unrounded value is what makes the two
regimes exactly continuous. -/
def tauonicInputMax : ℝ := muonMaxEjection
\mbox{}
/-- Maximum ejection velocity of the tauonic regime, at `v₂ = tauonicInputMax`. -/
def tauonMaxEjection : ℝ :=
  generalisedVelocity generalizedDimension tauonicInputMax massRatioTauon
\mbox{}
/-- The tauonic ejection velocity obeys the law: `vf² = v₂² · 3477³`. -/
theorem tauonMaxEjection\_\allowbreak{}sq :
    tauonMaxEjection \textasciicircum{} 2 = tauonicInputMax \textasciicircum{} 2 * massRatioTauon \textasciicircum{} 3 := by
  unfold tauonMaxEjection
  exact velocity\_\allowbreak{}law\_\allowbreak{}five \_\allowbreak{} \_\allowbreak{} massRatioTauon\_\allowbreak{}pos.le
\mbox{}
/-- `vf ≈ 4.2403 × 10¹² m/s` for the tauonic ceiling. -/
theorem tauonMaxEjection\_\allowbreak{}bounds :
    (4.2403 : ℝ) * 10 \textasciicircum{} 12 ≤ tauonMaxEjection ∧
      tauonMaxEjection ≤ (4.2404 : ℝ) * 10 \textasciicircum{} 12 := by
  have hb := muonMaxEjection\_\allowbreak{}bounds
  have hsq₁ : ((2.06820 : ℝ) * 10 \textasciicircum{} 7) \textasciicircum{} 2 ≤ tauonicInputMax \textasciicircum{} 2 := by
    rw [tauonicInputMax]
    exact pow\_\allowbreak{}le\_\allowbreak{}pow\_\allowbreak{}left₀ (show (0 : ℝ) ≤ (2.06820 : ℝ) * 10 \textasciicircum{} 7 by norm\_\allowbreak{}num) hb.1 2
  have hsq₂ : tauonicInputMax \textasciicircum{} 2 ≤ ((2.06821 : ℝ) * 10 \textasciicircum{} 7) \textasciicircum{} 2 := by
    rw [tauonicInputMax]
    exact pow\_\allowbreak{}le\_\allowbreak{}pow\_\allowbreak{}left₀ (by linarith [hb.1] : (0 : ℝ) ≤ muonMaxEjection) hb.2 2
  have hp₁ := mul\_\allowbreak{}le\_\allowbreak{}mul\_\allowbreak{}of\_\allowbreak{}nonneg\_\allowbreak{}right hsq₁
    (show (0 : ℝ) ≤ massRatioTauon \textasciicircum{} 3 by norm\_\allowbreak{}num [massRatioTauon])
  have hp₂ := mul\_\allowbreak{}le\_\allowbreak{}mul\_\allowbreak{}of\_\allowbreak{}nonneg\_\allowbreak{}right hsq₂
    (show (0 : ℝ) ≤ massRatioTauon \textasciicircum{} 3 by norm\_\allowbreak{}num [massRatioTauon])
  have h₁ : ((4.2403 : ℝ) * 10 \textasciicircum{} 12) \textasciicircum{} 2 ≤ tauonicInputMax \textasciicircum{} 2 * massRatioTauon \textasciicircum{} 3 := by
    refine le\_\allowbreak{}trans ?\_\allowbreak{} hp₁
    norm\_\allowbreak{}num [massRatioTauon]
  have h₂ : tauonicInputMax \textasciicircum{} 2 * massRatioTauon \textasciicircum{} 3 ≤ ((4.2404 : ℝ) * 10 \textasciicircum{} 12) \textasciicircum{} 2 := by
    refine le\_\allowbreak{}trans hp₂ ?\_\allowbreak{}
    norm\_\allowbreak{}num [massRatioTauon]
  have := generalisedVelocity\_\allowbreak{}bounds (n := generalizedDimension)
    (v₂ := tauonicInputMax) (r := massRatioTauon)
    (a := (4.2403 : ℝ) * 10 \textasciicircum{} 12) (b := (4.2404 : ℝ) * 10 \textasciicircum{} 12)
    (show (0 : ℝ) ≤ (4.2404 : ℝ) * 10 \textasciicircum{} 12 by norm\_\allowbreak{}num) h₁ h₂
  simpa [tauonMaxEjection] using this
\mbox{}
/-- The tauonic ceiling is superluminal: the framework's rejection of the
Lorentz factor is not optional decoration, it is what its own numbers require. -/
theorem tauonMaxEjection\_\allowbreak{}superluminal : speedOfLight \textless{} tauonMaxEjection := by
  have := tauonMaxEjection\_\allowbreak{}bounds.1
  unfold speedOfLight
  linarith
\mbox{}
/-- The tauonic input window is subluminal; only the output exceeds `c`. -/
theorem tauonicInputMax\_\allowbreak{}subluminal : tauonicInputMax \textless{} speedOfLight := by
  rw [tauonicInputMax]
  exact muonMaxEjection\_\allowbreak{}subluminal
\mbox{}
/-! \#\# The tauonic figures of the source text
\mbox{}
The source states `vf ≈ 1.34 × 10¹² m/s` at the tauonic ceiling, and
`vf ≈ 9.0 × 10¹⁰ m/s` at a median input of `500 000 km/h`.  The exact law gives
`4.2403 × 10¹²` and `2.8477 × 10¹⁰`.  In both cases the printed value is `√10`
too small: the source's own intermediate `vf² = 1.80 × 10²⁵ m²/s²` is correct,
but its square root is taken as `1.34 × 10¹²` instead of `4.24 × 10¹²`.  The
discrepancy is proved here rather than asserted, and both values are kept so
that a reader can see exactly where the two diverge.
-/
\mbox{}
/-- The tauonic ceiling ejection velocity as printed in the source text. -/
def sourceTauonMaxVelocity : ℝ := 1.34 * 10 \textasciicircum{} 12
\mbox{}
/-- The tauonic median ejection velocity as printed in the source text. -/
def sourceTauonMedianVelocity : ℝ := 9.0 * 10 \textasciicircum{} 10
\mbox{}
/-- The source figure, squared, falls short of the law by a factor of `10.01`:
the printed value is `√10` too small. -/
theorem source\_\allowbreak{}tauon\_\allowbreak{}max\_\allowbreak{}velocity\_\allowbreak{}underestimates :
    (10.01 : ℝ) ≤ tauonMaxEjection \textasciicircum{} 2 / sourceTauonMaxVelocity \textasciicircum{} 2 ∧
      tauonMaxEjection \textasciicircum{} 2 / sourceTauonMaxVelocity \textasciicircum{} 2 ≤ (10.02 : ℝ) := by
  obtain ⟨h₁, h₂⟩ := tauonMaxEjection\_\allowbreak{}bounds
  have hsq₁ : ((4.2403 : ℝ) * 10 \textasciicircum{} 12) \textasciicircum{} 2 ≤ tauonMaxEjection \textasciicircum{} 2 :=
    pow\_\allowbreak{}le\_\allowbreak{}pow\_\allowbreak{}left₀ (show (0 : ℝ) ≤ (4.2403 : ℝ) * 10 \textasciicircum{} 12 by norm\_\allowbreak{}num) h₁ 2
  have hsq₂ : tauonMaxEjection \textasciicircum{} 2 ≤ ((4.2404 : ℝ) * 10 \textasciicircum{} 12) \textasciicircum{} 2 :=
    pow\_\allowbreak{}le\_\allowbreak{}pow\_\allowbreak{}left₀ (by linarith : (0 : ℝ) ≤ tauonMaxEjection) h₂ 2
  unfold sourceTauonMaxVelocity
  constructor
  · rw [le\_\allowbreak{}div\_\allowbreak{}iff₀ (by norm\_\allowbreak{}num)]
    nlinarith
  · rw [div\_\allowbreak{}le\_\allowbreak{}iff₀ (by norm\_\allowbreak{}num)]
    nlinarith
\mbox{}
/-- Consequently the source figure underestimates the ejection velocity itself by
the factor `√10 ≈ 3.164`. -/
theorem source\_\allowbreak{}tauon\_\allowbreak{}max\_\allowbreak{}velocity\_\allowbreak{}factor :
    (3.164 : ℝ) ≤ tauonMaxEjection / sourceTauonMaxVelocity ∧
      tauonMaxEjection / sourceTauonMaxVelocity ≤ (3.165 : ℝ) := by
  obtain ⟨h₁, h₂⟩ := tauonMaxEjection\_\allowbreak{}bounds
  unfold sourceTauonMaxVelocity
  constructor
  · rw [le\_\allowbreak{}div\_\allowbreak{}iff₀ (by norm\_\allowbreak{}num)]
    nlinarith
  · rw [div\_\allowbreak{}le\_\allowbreak{}iff₀ (by norm\_\allowbreak{}num)]
    nlinarith
\mbox{}
/-- Tauonic ejection velocity at the entry velocity of the tauonic window. -/
def tauonEntryEjection : ℝ :=
  generalisedVelocity generalizedDimension collisionVelocityMuonMax massRatioTauon
\mbox{}
/-- `vf ≈ 1.4238 × 10⁹ m/s` at the tauonic entry velocity: the source's
`1.42 × 10⁹` is correct here, because this `vf²` happens to carry an even
decimal exponent. -/
theorem tauonEntryEjection\_\allowbreak{}bounds :
    (1.4237 : ℝ) * 10 \textasciicircum{} 9 ≤ tauonEntryEjection ∧
      tauonEntryEjection ≤ (1.4239 : ℝ) * 10 \textasciicircum{} 9 := by
  have h₁ : ((1.4237 : ℝ) * 10 \textasciicircum{} 9) \textasciicircum{} 2 ≤
      collisionVelocityMuonMax \textasciicircum{} 2 * massRatioTauon \textasciicircum{} 3 := by
    rw [collisionVelocityMuonMax\_\allowbreak{}eq]
    norm\_\allowbreak{}num [massRatioTauon]
  have h₂ : collisionVelocityMuonMax \textasciicircum{} 2 * massRatioTauon \textasciicircum{} 3 ≤
      ((1.4239 : ℝ) * 10 \textasciicircum{} 9) \textasciicircum{} 2 := by
    rw [collisionVelocityMuonMax\_\allowbreak{}eq]
    norm\_\allowbreak{}num [massRatioTauon]
  have := generalisedVelocity\_\allowbreak{}bounds (n := generalizedDimension)
    (v₂ := collisionVelocityMuonMax) (r := massRatioTauon)
    (a := (1.4237 : ℝ) * 10 \textasciicircum{} 9) (b := (1.4239 : ℝ) * 10 \textasciicircum{} 9)
    (show (0 : ℝ) ≤ (1.4239 : ℝ) * 10 \textasciicircum{} 9 by norm\_\allowbreak{}num) h₁ h₂
  simpa [tauonEntryEjection] using this
\mbox{}
/-- Tauonic ejection velocity at a median input of `500 000 km/h`. -/
def tauonMedianEjection : ℝ :=
  generalisedVelocity generalizedDimension (kmhToMs 500000) massRatioTauon
\mbox{}
/-- `vf ≈ 2.8477 × 10¹⁰ m/s` at the median input — not the `9.0 × 10¹⁰` printed
in the source, which carries the same `√10` slip. -/
theorem tauonMedianEjection\_\allowbreak{}bounds :
    (2.8475 : ℝ) * 10 \textasciicircum{} 10 ≤ tauonMedianEjection ∧
      tauonMedianEjection ≤ (2.8476 : ℝ) * 10 \textasciicircum{} 10 := by
  have hvm : kmhToMs 500000 = 1250000 / 9 := by
    unfold kmhToMs
    norm\_\allowbreak{}num
  have h₁ : ((2.8475 : ℝ) * 10 \textasciicircum{} 10) \textasciicircum{} 2 ≤
      kmhToMs 500000 \textasciicircum{} 2 * massRatioTauon \textasciicircum{} 3 := by
    rw [hvm]
    norm\_\allowbreak{}num [massRatioTauon]
  have h₂ : kmhToMs 500000 \textasciicircum{} 2 * massRatioTauon \textasciicircum{} 3 ≤
      ((2.8476 : ℝ) * 10 \textasciicircum{} 10) \textasciicircum{} 2 := by
    rw [hvm]
    norm\_\allowbreak{}num [massRatioTauon]
  have := generalisedVelocity\_\allowbreak{}bounds (n := generalizedDimension)
    (v₂ := kmhToMs 500000) (r := massRatioTauon)
    (a := (2.8475 : ℝ) * 10 \textasciicircum{} 10) (b := (2.8476 : ℝ) * 10 \textasciicircum{} 10)
    (show (0 : ℝ) ≤ (2.8476 : ℝ) * 10 \textasciicircum{} 10 by norm\_\allowbreak{}num) h₁ h₂
  simpa [tauonMedianEjection] using this
\mbox{}
end
\mbox{}
end TheoreticalFramework
\mbox{}
}

\clearpage


\subsection*{\texttt{TheoreticalFramework/Regimes.lean}}
\vspace{2pt}
{\ttfamily\scriptsize\setlength{\parindent}{0pt}\setlength{\baselineskip}{7.6pt}\setlength{\parskip}{0pt}\raggedright\obeyspaces\obeylines
/-
\# TheoreticalFramework.Regimes
\mbox{}
The two nested propulsive regimes, their continuity at the domain boundary, the
structural algebra of the amplification factor, and an exact treatment of what
the framework does *not* fix: thrust.
\mbox{}
Structure of the argument.
\mbox{}
* A `PropulsiveRegime` is a mass ratio together with the window of admissible
  collision velocities `v₂` and the law-induced output.  Both regimes are
  instances, so everything proved about `PropulsiveRegime` applies to both.
* The muonic window is `]0, 25 000 km/h]`; the tauonic window runs from
  `25 000 km/h` to the muonic regime's *output* ceiling, `≈ 20 682 km/s`.
  Hence the windows tile one interval (`regimes\_\allowbreak{}tile`), the tauonic window is
  nested inside the muonic output range (`tauonic\_\allowbreak{}window\_\allowbreak{}nested\_\allowbreak{}in\_\allowbreak{}output`), and
  the output ceiling of the first regime is the input ceiling of the second
  (`chain\_\allowbreak{}closes`).  That is the exact content of the continuity claim.
* The amplification `r ↦ r\textasciicircum{}(n−2)` is a multiplicative character
  (`amplification\_\allowbreak{}mul`), so the tauonic factor factors through the muonic one.
* Thrust is *not* determined by the geometry.  The kinematics fix the exhaust
  velocity and the thrust-to-power ratio `T/P = 2/v\_\allowbreak{}e` exactly; the absolute
  thrust is free (`thrust\_\allowbreak{}not\_\allowbreak{}determined`).  What remains is an engineering
  question, recorded as `ThrustOptimizationConjecture` and left unproven, which
  is what the framework claims and no more.
-/
import TheoreticalFramework.Velocity
\mbox{}
set\_\allowbreak{}option linter.unusedVariables false
\mbox{}
namespace TheoreticalFramework
\mbox{}
noncomputable section
\mbox{}
/-! \#\# Amplification as a multiplicative character -/
\mbox{}
/-- The amplification factor is multiplicative in the mass ratio: chaining two
mass scales multiplies the amplification.  This is what makes the regimes
composable rather than merely analogous. -/
theorem amplification\_\allowbreak{}mul (r₁ r₂ : ℝ) :
    (r₁ * r₂) \textasciicircum{} velocityExponent = r₁ \textasciicircum{} velocityExponent * r₂ \textasciicircum{} velocityExponent :=
  mul\_\allowbreak{}pow r₁ r₂ \_\allowbreak{}
\mbox{}
/-- The amplification is strictly increasing in the mass ratio on the positive
axis: heavier collision partners amplify more. -/
theorem amplification\_\allowbreak{}strictMono :
    StrictMonoOn (fun r : ℝ =\textgreater{} r \textasciicircum{} velocityExponent) (Set.Ioi 0) := by
  intro r₁ hr₁ r₂ hr₂ hlt
  exact pow\_\allowbreak{}lt\_\allowbreak{}pow\_\allowbreak{}left₀ hlt (le\_\allowbreak{}of\_\allowbreak{}lt hr₁) (by decide)
\mbox{}
/-- The tauonic amplification is the muonic one times the amplification of the
residual ratio `m\_\allowbreak{}τ/m\_\allowbreak{}μ`: `3477³ = 207³ · (3477/207)³`. -/
theorem tauon\_\allowbreak{}amplification\_\allowbreak{}factors :
    tauonAmplification = muonAmplification * (massRatioTauon / massRatioMuon) \textasciicircum{} velocityExponent := by
  unfold tauonAmplification muonAmplification
  rw [div\_\allowbreak{}pow, mul\_\allowbreak{}div\_\allowbreak{}cancel₀ \_\allowbreak{} (pow\_\allowbreak{}ne\_\allowbreak{}zero \_\allowbreak{} (ne\_\allowbreak{}of\_\allowbreak{}gt massRatioMuon\_\allowbreak{}pos))]
\mbox{}
/-! \#\# The ejection velocity never falls below the collision velocity -/
\mbox{}
/-- For a mass ratio at least one, the ejection velocity is at least the
collision velocity.  This is what lets the tauonic window start where the muonic
window ends. -/
theorem generalisedVelocity\_\allowbreak{}ge\_\allowbreak{}input \{n : ℕ\} \{v₂ r : ℝ\} (hr : 1 ≤ r) :
    v₂ ≤ generalisedVelocity n v₂ r := by
  have hone : (1 : ℝ) ≤ r \textasciicircum{} (n - 2) := one\_\allowbreak{}le\_\allowbreak{}pow₀ hr
  have hsq : v₂ \textasciicircum{} 2 ≤ v₂ \textasciicircum{} 2 * r \textasciicircum{} (n - 2) := le\_\allowbreak{}mul\_\allowbreak{}of\_\allowbreak{}one\_\allowbreak{}le\_\allowbreak{}right (sq\_\allowbreak{}nonneg v₂) hone
  exact Real.le\_\allowbreak{}sqrt\_\allowbreak{}of\_\allowbreak{}sq\_\allowbreak{}le hsq
\mbox{}
theorem one\_\allowbreak{}le\_\allowbreak{}massRatioMuon : (1 : ℝ) ≤ massRatioMuon := by
  norm\_\allowbreak{}num [massRatioMuon]
\mbox{}
theorem one\_\allowbreak{}le\_\allowbreak{}massRatioTauon : (1 : ℝ) ≤ massRatioTauon := by
  norm\_\allowbreak{}num [massRatioTauon]
\mbox{}
/-! \#\# Regimes -/
\mbox{}
/-- A propulsive regime of the framework: a charged-lepton mass ratio, the
window of admissible collision velocities, and the law-induced quantities
derived from them. -/
structure PropulsiveRegime where
  /-- Display name of the regime. -/
  name : String
  /-- Colliding particle, for the record. -/
  particle : String
  /-- Mass ratio `m\_\allowbreak{}ℓ/m\_\allowbreak{}e` of the colliding leptons. -/
  massRatio : ℝ
  massRatio\_\allowbreak{}pos : 0 \textless{} massRatio
  /-- Lower end of the admissible input window, in m/s. -/
  inputLower : ℝ
  /-- Upper end of the admissible input window, in m/s. -/
  inputUpper : ℝ
  input\_\allowbreak{}window : inputLower ≤ inputUpper
\mbox{}
namespace PropulsiveRegime
\mbox{}
/-- Amplification factor of the regime, `(m\_\allowbreak{}ℓ/m\_\allowbreak{}e)\textasciicircum{}(n−2)`. -/
def amplification (R : PropulsiveRegime) : ℝ := R.massRatio \textasciicircum{} velocityExponent
\mbox{}
/-- Maximum ejection velocity of the regime, at the top of its input window. -/
def maxOutput (R : PropulsiveRegime) : ℝ :=
  generalisedVelocity generalizedDimension R.inputUpper R.massRatio
\mbox{}
/-- The input window as a set. -/
def window (R : PropulsiveRegime) : Set ℝ := Set.Icc R.inputLower R.inputUpper
\mbox{}
/-- The regime's maximum output obeys the velocity law. -/
theorem maxOutput\_\allowbreak{}sq (R : PropulsiveRegime) :
    R.maxOutput \textasciicircum{} 2 = R.inputUpper \textasciicircum{} 2 * R.massRatio \textasciicircum{} 3 := by
  unfold maxOutput
  exact velocity\_\allowbreak{}law\_\allowbreak{}five \_\allowbreak{} \_\allowbreak{} (le\_\allowbreak{}of\_\allowbreak{}lt R.massRatio\_\allowbreak{}pos)
\mbox{}
/-- The maximum output is at least the input ceiling. -/
theorem inputUpper\_\allowbreak{}le\_\allowbreak{}maxOutput (R : PropulsiveRegime) (h : 1 ≤ R.massRatio) :
    R.inputUpper ≤ R.maxOutput :=
  generalisedVelocity\_\allowbreak{}ge\_\allowbreak{}input h
\mbox{}
/-- Within a regime the velocity law is continuous in the collision velocity. -/
theorem continuous\_\allowbreak{}maxOutput (R : PropulsiveRegime) :
    Continuous fun v₂ : ℝ =\textgreater{} generalisedVelocity generalizedDimension v₂ R.massRatio := by
  unfold generalisedVelocity
  exact Real.continuous\_\allowbreak{}sqrt.comp (by fun\_\allowbreak{}prop)
\mbox{}
/-- Within a regime the velocity law is continuous in the mass ratio. -/
theorem continuous\_\allowbreak{}in\_\allowbreak{}massRatio (R : PropulsiveRegime) (v₂ : ℝ) :
    Continuous fun r : ℝ =\textgreater{} generalisedVelocity generalizedDimension v₂ r := by
  unfold generalisedVelocity
  refine Real.continuous\_\allowbreak{}sqrt.comp ?\_\allowbreak{}
  exact continuous\_\allowbreak{}const.mul (continuous\_\allowbreak{}id.pow \_\allowbreak{})
\mbox{}
end PropulsiveRegime
\mbox{}
/-- **The muonic regime.**  Muon–muon collisions, mass ratio `207`, collision
velocities up to `25 000 km/h`; replaces chemical propellants. -/
def muonicRegime : PropulsiveRegime where
  name := "muonic"
  particle := "μ⁻/μ⁺"
  massRatio := massRatioMuon
  massRatio\_\allowbreak{}pos := massRatioMuon\_\allowbreak{}pos
  inputLower := 0
  inputUpper := collisionVelocityMuonMax
  input\_\allowbreak{}window := by
    have h := collisionVelocityMuonMax\_\allowbreak{}bounds.1
    linarith
\mbox{}
/-- **The tauonic regime.**  Tauon–tauon collisions, mass ratio `3477`, input
window from the muonic ceiling up to the muonic regime's output ceiling. -/
def tauonicRegime : PropulsiveRegime where
  name := "tauonic"
  particle := "τ⁻/τ⁺"
  massRatio := massRatioTauon
  massRatio\_\allowbreak{}pos := massRatioTauon\_\allowbreak{}pos
  inputLower := collisionVelocityMuonMax
  inputUpper := tauonicInputMax
  input\_\allowbreak{}window := by
    have h := generalisedVelocity\_\allowbreak{}ge\_\allowbreak{}input
      (n := generalizedDimension) (v₂ := collisionVelocityMuonMax) (r := massRatioMuon)
      one\_\allowbreak{}le\_\allowbreak{}massRatioMuon
    simpa [tauonicInputMax, muonMaxEjection] using h
\mbox{}
/-- The muonic amplification, through the structure. -/
theorem muonicRegime\_\allowbreak{}amplification :
    muonicRegime.amplification = muonAmplification := by
  unfold PropulsiveRegime.amplification muonicRegime muonAmplification
  rfl
\mbox{}
/-- The tauonic amplification, through the structure. -/
theorem tauonicRegime\_\allowbreak{}amplification :
    tauonicRegime.amplification = tauonAmplification := by
  unfold PropulsiveRegime.amplification tauonicRegime tauonAmplification
  rfl
\mbox{}
/-- The regime hierarchy is strict: the tauonic amplification exceeds the
muonic one by `≈ 4.74 × 10³`. -/
theorem regimes\_\allowbreak{}strictly\_\allowbreak{}ordered :
    muonicRegime.amplification \textless{} tauonicRegime.amplification := by
  rw [muonicRegime\_\allowbreak{}amplification, tauonicRegime\_\allowbreak{}amplification,
    muonAmplification\_\allowbreak{}eq, tauonAmplification\_\allowbreak{}eq]
  norm\_\allowbreak{}num
\mbox{}
/-- The muonic regime's output ceiling, through the structure. -/
theorem muonicRegime\_\allowbreak{}maxOutput : muonicRegime.maxOutput = muonMaxEjection := by
  unfold PropulsiveRegime.maxOutput muonicRegime muonMaxEjection
  rfl
\mbox{}
/-- The tauonic regime's output ceiling, through the structure. -/
theorem tauonicRegime\_\allowbreak{}maxOutput : tauonicRegime.maxOutput = tauonMaxEjection := by
  unfold PropulsiveRegime.maxOutput tauonicRegime tauonMaxEjection
  rfl
\mbox{}
/-! \#\# Continuity at the domain boundary -/
\mbox{}
/-- **Adjacency.**  The tauonic window begins exactly where the muonic window
ends: `25 000 km/h`, the muonic collision ceiling. -/
theorem windows\_\allowbreak{}adjacent : muonicRegime.inputUpper = tauonicRegime.inputLower := rfl
\mbox{}
/-- **The chain closes.**  The upper end of the tauonic *input* window is the
muonic regime's maximum *output*.  This is the boundary the framework calls
`v₂ = 20 700 km/s`, and taking the unrounded value is what makes the match
exact rather than approximate. -/
theorem chain\_\allowbreak{}closes : tauonicRegime.inputUpper = muonicRegime.maxOutput := rfl
\mbox{}
/-- **Tiling.**  The two input windows cover one interval with no gap and no
overlap. -/
theorem regimes\_\allowbreak{}tile :
    muonicRegime.window ∪ tauonicRegime.window
      = Set.Icc (0 : ℝ) tauonicRegime.inputUpper := by
  ext v
  simp only [PropulsiveRegime.window, muonicRegime, tauonicRegime, Set.mem\_\allowbreak{}union,
    Set.mem\_\allowbreak{}Icc]
  constructor
  · rintro (⟨h₀, h₁⟩ | ⟨h₁, h₂⟩)
    · exact ⟨h₀, le\_\allowbreak{}trans h₁ (by
        have := generalisedVelocity\_\allowbreak{}ge\_\allowbreak{}input (n := generalizedDimension)
          (v₂ := collisionVelocityMuonMax) (r := massRatioMuon) one\_\allowbreak{}le\_\allowbreak{}massRatioMuon
        simpa [tauonicInputMax, muonMaxEjection] using this)⟩
    · exact ⟨le\_\allowbreak{}trans (by
        have := collisionVelocityMuonMax\_\allowbreak{}bounds
        linarith [this.1]) h₁, h₂⟩
  · rintro ⟨h₀, h₂⟩
    rcases le\_\allowbreak{}total v collisionVelocityMuonMax with h | h
    · exact Or.inl ⟨h₀, h⟩
    · exact Or.inr ⟨h, h₂⟩
\mbox{}
/-- **Nesting.**  The tauonic window lies inside the muonic regime's output
range: the second regime's inputs are the first regime's reachable outputs. -/
theorem tauonic\_\allowbreak{}window\_\allowbreak{}nested\_\allowbreak{}in\_\allowbreak{}output :
    tauonicRegime.window ⊆ Set.Ioc (0 : ℝ) muonicRegime.maxOutput := by
  intro v hv
  have h := regimes\_\allowbreak{}tile
  simp only [PropulsiveRegime.window, tauonicRegime, Set.mem\_\allowbreak{}Icc] at hv
  refine ⟨?\_\allowbreak{}, ?\_\allowbreak{}⟩
  · have := collisionVelocityMuonMax\_\allowbreak{}bounds
    linarith [hv.1]
  · rw [tauonicInputMax] at hv
    exact hv.2
\mbox{}
/-- **The domain boundary.**  The boundary of the source's §8.2, `v₂ = 20 700 km/s`,
is the tauonic input ceiling, and it is the muonic output ceiling unrounded. -/
def domainBoundary : ℝ := tauonicRegime.inputUpper
\mbox{}
/-- The boundary is the muonic regime's maximum output. -/
theorem domainBoundary\_\allowbreak{}eq\_\allowbreak{}muonic\_\allowbreak{}output : domainBoundary = muonicRegime.maxOutput := rfl
\mbox{}
/-- The boundary is the muonic ejection ceiling, hence `≈ 20 700 km/s`. -/
theorem domainBoundary\_\allowbreak{}km\_\allowbreak{}per\_\allowbreak{}s :
    (2.068 : ℝ) * 10 \textasciicircum{} 4 ≤ domainBoundary / 1000 ∧
      domainBoundary / 1000 ≤ (2.07 : ℝ) * 10 \textasciicircum{} 4 := by
  simpa [domainBoundary, tauonicRegime, tauonicInputMax] using muonMaxEjection\_\allowbreak{}km\_\allowbreak{}per\_\allowbreak{}s
\mbox{}
/-- The boundary itself is subluminal; it is the tauonic *output* that exceeds
`c`, not the input ceiling that defines the hand-off. -/
theorem domainBoundary\_\allowbreak{}subluminal : domainBoundary \textless{} speedOfLight := by
  simpa [domainBoundary, tauonicRegime, tauonicInputMax] using muonMaxEjection\_\allowbreak{}subluminal
\mbox{}
/-- **Continuity of the two-regime structure**, collected in one statement: the
windows are adjacent, the chain closes on the muonic output ceiling, they tile
an interval, and the law is continuous in the collision velocity throughout. -/
theorem regime\_\allowbreak{}continuity :
    muonicRegime.inputUpper = tauonicRegime.inputLower ∧
      tauonicRegime.inputUpper = muonicRegime.maxOutput ∧
      muonicRegime.window ∪ tauonicRegime.window
        = Set.Icc (0 : ℝ) tauonicRegime.inputUpper ∧
      Continuous (fun v₂ : ℝ =\textgreater{}
        generalisedVelocity generalizedDimension v₂ muonicRegime.massRatio) ∧
      Continuous (fun v₂ : ℝ =\textgreater{}
        generalisedVelocity generalizedDimension v₂ tauonicRegime.massRatio) :=
  ⟨windows\_\allowbreak{}adjacent, chain\_\allowbreak{}closes, regimes\_\allowbreak{}tile,
    muonicRegime.continuous\_\allowbreak{}maxOutput, tauonicRegime.continuous\_\allowbreak{}maxOutput⟩
\mbox{}
/-! \#\# Thrust: what the framework fixes and what it does not
\mbox{}
The kinematics fix the exhaust velocity.  Momentum and energy balance then fix
the thrust-to-power ratio exactly, `T/P = 2/v\_\allowbreak{}e`, and nothing else: the absolute
thrust is set by mass flow and power, which are engineering parameters.  This is
the precise sense in which "empirical thrust optimisation" is outside the
framework, and it is proved rather than asserted.
-/
\mbox{}
/-- A thruster design obeying the framework's momentum and energy balance. -/
structure ThrusterDesign where
  /-- Exhaust velocity of the propulsive particles, m/s. -/
  exhaustVelocity : ℝ
  exhaustVelocity\_\allowbreak{}pos : 0 \textless{} exhaustVelocity
  /-- Propellant mass flow, kg/s. -/
  massFlow : ℝ
  massFlow\_\allowbreak{}pos : 0 \textless{} massFlow
  /-- Thrust, N: momentum flux `ṁ·v\_\allowbreak{}e`. -/
  thrust : ℝ
  thrust\_\allowbreak{}eq : thrust = massFlow * exhaustVelocity
  /-- Jet power, W: kinetic energy flux `½·ṁ·v\_\allowbreak{}e²`. -/
  jetPower : ℝ
  power\_\allowbreak{}eq : jetPower = massFlow * exhaustVelocity \textasciicircum{} 2 / 2
\mbox{}
/-- Vacuum exhaust velocity of the Raptor benchmark. -/
def raptorExhaustVelocity : ℝ := raptorSpecificImpulse * standardGravity
\mbox{}
/-- **The thrust–power tradeoff.**  For any design, `T/P = 2/v\_\allowbreak{}e` exactly.  High
specific impulse costs thrust per watt, without exception. -/
theorem thrust\_\allowbreak{}to\_\allowbreak{}power (d : ThrusterDesign) :
    d.thrust / d.jetPower = 2 / d.exhaustVelocity := by
  rw [d.thrust\_\allowbreak{}eq, d.power\_\allowbreak{}eq]
  field\_\allowbreak{}simp [ne\_\allowbreak{}of\_\allowbreak{}gt d.massFlow\_\allowbreak{}pos, ne\_\allowbreak{}of\_\allowbreak{}gt d.exhaustVelocity\_\allowbreak{}pos]
\mbox{}
/-- Thrust in terms of jet power and exhaust velocity. -/
theorem thrust\_\allowbreak{}eq\_\allowbreak{}power (d : ThrusterDesign) :
    d.thrust = 2 * d.jetPower / d.exhaustVelocity := by
  rw [d.thrust\_\allowbreak{}eq, d.power\_\allowbreak{}eq]
  field\_\allowbreak{}simp [ne\_\allowbreak{}of\_\allowbreak{}gt d.exhaustVelocity\_\allowbreak{}pos]
\mbox{}
/-- At the muonic ceiling, thrust per watt is `≈ 9.67 × 10⁻⁸ N/W`. -/
theorem muonic\_\allowbreak{}thrust\_\allowbreak{}to\_\allowbreak{}power\_\allowbreak{}bounds (d : ThrusterDesign)
    (h : d.exhaustVelocity = muonMaxEjection) :
    (96 : ℝ) / 10 \textasciicircum{} 9 ≤ d.thrust / d.jetPower ∧
      d.thrust / d.jetPower ≤ (97 : ℝ) / 10 \textasciicircum{} 9 := by
  rw [thrust\_\allowbreak{}to\_\allowbreak{}power, h]
  obtain ⟨h₁, h₂⟩ := muonMaxEjection\_\allowbreak{}bounds
  constructor
  · rw [le\_\allowbreak{}div\_\allowbreak{}iff₀ (by linarith)]
    nlinarith
  · rw [div\_\allowbreak{}le\_\allowbreak{}iff₀ (by linarith)]
    nlinarith
\mbox{}
/-- The muonic exhaust velocity is `≈ 5.8 × 10³` times the Raptor exhaust
velocity, so its thrust per watt is that many times *worse*: the specific-impulse
gain and the thrust-per-watt loss are the same number. -/
theorem muonic\_\allowbreak{}thrust\_\allowbreak{}to\_\allowbreak{}power\_\allowbreak{}shortfall :
    (5.80 : ℝ) * 10 \textasciicircum{} 3 ≤ muonMaxEjection / raptorExhaustVelocity ∧
      muonMaxEjection / raptorExhaustVelocity ≤ (5.82 : ℝ) * 10 \textasciicircum{} 3 := by
  have hrv : raptorExhaustVelocity = 363 * 9.80665 := by
    unfold raptorExhaustVelocity raptorSpecificImpulse standardGravity
    norm\_\allowbreak{}num
  obtain ⟨h₁, h₂⟩ := muonMaxEjection\_\allowbreak{}bounds
  constructor
  · rw [hrv, le\_\allowbreak{}div\_\allowbreak{}iff₀ (by norm\_\allowbreak{}num)]
    nlinarith
  · rw [hrv, div\_\allowbreak{}le\_\allowbreak{}iff₀ (by norm\_\allowbreak{}num)]
    nlinarith
\mbox{}
/-- **Thrust is not determined by the geometry.**  At the exhaust velocity the
leptonic kinematics fix, any positive thrust is attainable by a suitable design:
thrust is bought with mass flow and power, not with the velocity law. -/
theorem thrust\_\allowbreak{}not\_\allowbreak{}determined (T : ℝ) (hT : 0 \textless{} T) :
    ∃ d : ThrusterDesign, d.exhaustVelocity = muonMaxEjection ∧ d.thrust = T := by
  have hv : 0 \textless{} muonMaxEjection := by
    have := muonMaxEjection\_\allowbreak{}bounds.1
    linarith
  refine ⟨\{ exhaustVelocity := muonMaxEjection, exhaustVelocity\_\allowbreak{}pos := hv,
            massFlow := T / muonMaxEjection, massFlow\_\allowbreak{}pos := div\_\allowbreak{}pos hT hv,
            thrust := T, thrust\_\allowbreak{}eq := ?\_\allowbreak{}, jetPower := T * muonMaxEjection / 2,
            power\_\allowbreak{}eq := ?\_\allowbreak{} \}, rfl, rfl⟩
  · field\_\allowbreak{}simp
  · field\_\allowbreak{}simp
\mbox{}
/-- The engineering question the framework leaves open: at the exhaust velocity
the geometry fixes, is there a design reaching a thrust target `T` within a
power budget `P`?  Nothing in the geometric development decides it — see
`thrust\_\allowbreak{}not\_\allowbreak{}determined` — so it is recorded, not proved. -/
def ThrustOptimizationConjecture (T P : ℝ) : Prop :=
  ∃ d : ThrusterDesign, d.exhaustVelocity = muonMaxEjection ∧ d.thrust = T ∧ d.jetPower ≤ P
\mbox{}
/-- The conjecture is not vacuous: a large enough power budget always suffices,
because the thrust–power relation is a hyperbola and nothing else constrains it. -/
theorem thrust\_\allowbreak{}optimization\_\allowbreak{}satisfiable (T : ℝ) (hT : 0 \textless{} T) :
    ThrustOptimizationConjecture T (T * muonMaxEjection / 2 + 1) := by
  obtain ⟨d, hv, hth⟩ := thrust\_\allowbreak{}not\_\allowbreak{}determined T hT
  refine ⟨d, hv, hth, ?\_\allowbreak{}⟩
  have hvpos : 0 \textless{} muonMaxEjection := by
    have := muonMaxEjection\_\allowbreak{}bounds.1
    linarith
  have hmv : d.massFlow * d.exhaustVelocity = T := by
    rw [← d.thrust\_\allowbreak{}eq, hth]
  rw [hv] at hmv
  have hmf : d.massFlow = T / muonMaxEjection := by
    rw [eq\_\allowbreak{}div\_\allowbreak{}iff (ne\_\allowbreak{}of\_\allowbreak{}gt hvpos)]
    linarith
  have hp := d.power\_\allowbreak{}eq
  rw [hv, hmf] at hp
  rw [hp]
  have hnorm : T / muonMaxEjection * muonMaxEjection \textasciicircum{} 2 / 2 = T * muonMaxEjection / 2 := by
    field\_\allowbreak{}simp
  rw [hnorm]
  linarith
\mbox{}
/-- Consequently the thrust question is *empirical*, not geometric: the
framework fixes `v\_\allowbreak{}e` and hence `T/P`, and leaves `(T, P)` free along that
hyperbola.  This is the formal statement of "empirical thrust optimisations
remain unproven". -/
theorem thrust\_\allowbreak{}optimization\_\allowbreak{}is\_\allowbreak{}empirical (T P : ℝ) (hT : 0 \textless{} T) (hP : 0 \textless{} P)
    (hpower : T ≤ 2 * P / muonMaxEjection) :
    ThrustOptimizationConjecture T P := by
  have hv : 0 \textless{} muonMaxEjection := by
    have := muonMaxEjection\_\allowbreak{}bounds.1
    linarith
  refine ⟨\{ exhaustVelocity := muonMaxEjection, exhaustVelocity\_\allowbreak{}pos := hv,
            massFlow := T / muonMaxEjection, massFlow\_\allowbreak{}pos := div\_\allowbreak{}pos hT hv,
            thrust := T, thrust\_\allowbreak{}eq := ?\_\allowbreak{}, jetPower := T * muonMaxEjection / 2,
            power\_\allowbreak{}eq := ?\_\allowbreak{} \}, rfl, rfl, ?\_\allowbreak{}⟩
  · field\_\allowbreak{}simp
  · field\_\allowbreak{}simp
  · have hp : T * muonMaxEjection ≤ 2 * P := by
      have h' := hpower
      rw [le\_\allowbreak{}div\_\allowbreak{}iff₀ hv] at h'
      linarith
    show T * muonMaxEjection / 2 ≤ P
    rw [div\_\allowbreak{}le\_\allowbreak{}iff₀ (by norm\_\allowbreak{}num : (0 : ℝ) \textless{} 2)]
    linarith
\mbox{}
end
\mbox{}
end TheoreticalFramework
\mbox{}
}

\clearpage

\end{document}
